\documentclass[11pt]{ctexart}
\usepackage[a4paper,margin=1in]{geometry}
\usepackage{graphicx}
\usepackage{xcolor}
\usepackage{hyperref}
\definecolor{refgray}{gray}{0.45}
\title{\textbf{Market Views｜全球投行叙事汇编}\\[0.4em]{\large AI变现进入验证期，全球资金再配置与消费K型分化加剧}}
\author{KC桌面 自动生成}
\date{260730}
\begin{document}
\maketitle
\tableofcontents
\newpage
\section{一页摘要}
\begin{itemize}
  \item 全球权益资金流入创历史新高，但美国流入放缓至3月来最低，欧洲资金回流迹象显现。
  \item 中国消费呈K型分化：高端白酒与奢侈品韧性明显，但整体零售增速仅0.2\%，大众消费承压。
  \item AI产业从模型竞赛转向生态博弈，企业采购重心从‘选模型’变为‘选生态’，边缘推理与SaaS货币化成为关键战场。
  \item 存储需求从HBM扩散至传统DRAM/NAND，长期协议约束力增强，价格波动性结构性下降。
  \item 半导体设备需求强劲，Lam Research业绩超预期，毛利率创20年新高；Seagate近线HDD产能锁定至2028年。
  \item 电网投资持续高景气，变压器供需缺口预计2028年才收窄，中国供应商凭借短交货周期填补缺口。
  \item Meta和微软的广告与Copilot业务已产生可量化AI回报，但消费者AI粘性差异与主题乐园客流增速放缓构成不确定性。
  \item 亚洲市场分化加剧：中国智慧城市贡献20\%增长，而其余地区增速放缓；科技板块去杠杆接近尾声。
\end{itemize}
\section{全球投行叙事汇编}
今日纳入 63 篇投行/券商研究，按 9 个市场、资产与产业主题系统整理。
\newpage
\subsection{宏观与利率：全球资金再配置与消费分化}
\textbf{全球权益资金流入创纪录但结构生变，美国流入放缓而欧洲回流；中国消费呈K型分化，高端韧性但大众承压；猪周期因效率提升而延迟；AI存储需求扩散至传统DRAM/NAND，长期协议改变盈利稳定性。}
\subsubsection*{主流共识}
\begin{itemize}
  \item 全球权益资金流入创历史新高，但美国流入节奏放缓至3月以来最低，欧洲出现资金回流迹象。
  \item 中国消费市场呈K型分化：高端白酒、奢侈品韧性明显，但整体零售增速仅0.2\%，大众消费恢复面临阻力。
  \item AI对存储的需求从HBM扩散至传统DRAM和NAND，2027年DRAM位元消耗量增速预计从22\%升至36\%。
\end{itemize}
\subsubsection*{分歧与条件差异}
\begin{itemize}
  \item 猪周期拐点时间：MS认为产能出清慢于预期，价格反弹推迟至2027年三季度；而市场此前普遍预期2026年下半年见底。
  \item 美国消费者韧性：Bernstein认为消费者系统性转向价值导向，低收入群体承压；但整体失业率低位支撑消费支出。
  \item 奔驰中国业务前景：BARC和Bernstein均指出中国利润贡献下降，但管理层重申长期承诺，市场对减值影响定价是否充分存在分歧。
\end{itemize}
\subsubsection*{机构观点}
\paragraph{BARC} 梅赛德斯-奔驰信用质量在评级区间偏弱，市场未充分定价中国利润下降、成本上升与激进股东回报的叠加影响。H1自由现金流同比降41\%至11亿欧元，但工业净流动性仍达304亿欧元；管理层预计H2乘用车利润率弱于H1，全年利润率落在3\%-5\%区间下沿。
{\small\color{refgray}数据：H1自由现金流同比降41\%至11亿欧元\par}
{\small\color{refgray}数据：工业净流动性304亿欧元\par}
{\small\color{refgray}数据：H1股东回报50亿欧元\par}
{\small\color{refgray}数据：中国合资公司减值7.96亿欧元\par}
{\small\color{refgray}边际变化：此前市场关注流动性充裕，现转向现金流生成能力被结构性因素侵蚀，且中国业务利润预期显著下调。\par}
\paragraph{BARC} 7月全球权益净流入1290亿美元，年初至今累计6590亿美元创纪录，但美国流入降至3月以来最低，欧洲出现回流。对冲产品和CTA多头重新累积，整体持仓接近历史高位，共同基金现金水平处于十年底部，市场应对冲击的缓冲有限。
{\small\color{refgray}数据：7月全球权益净流入1290亿美元\par}
{\small\color{refgray}数据：年初至今累计流入6590亿美元创纪录\par}
{\small\color{refgray}数据：美国流入降至3月以来最低\par}
{\small\color{refgray}数据：共同基金现金水平处于十年底部\par}
{\small\color{refgray}边际变化：此前资金持续流入美国，现结构转变：美国放缓、欧洲回流，且零售情绪降温，增量资金转向机构被动配置。\par}
\paragraph{Bernstein} 中国Q2 GDP增速从5.0\%降至4.3\%，零售仅增0.2\%，消费呈K型分化。高端白酒批价上涨13\%-30\%，茅台年内第二次提价；啤酒动销转正（+2\%）。运动服饰线上增长10\%但线下疲弱。酒店高端与经济型RevPAR强劲，中间段承压。
{\small\color{refgray}数据：Q2 GDP增速4.3\%（Q1为5.0\%）\par}
{\small\color{refgray}数据：零售增速0.2\%\par}
{\small\color{refgray}数据：飞天茅台批价较底部涨13\%\par}
{\small\color{refgray}数据：啤酒Q2销售同比+2\%（Q1为-1\%）\par}
{\small\color{refgray}边际变化：此前市场关注消费信心修复，现发现修复呈K型：高端消费韧性超预期，但大众可选消费恢复缓慢。\par}
\paragraph{Bernstein} Equinix Q2业绩超预期，将年度AFFO每股增速从5\%-9\%上调至9\%-12\%，资本开支从30-40亿提升至50-70亿美元。约80\%投入25个核心市场，开发回报约25\%。AI需求已嵌入核心业务：前十大AI模型提供商中八家使用其网络，新增9700个互联节点创纪录。
{\small\color{refgray}数据：AFFO每股增速指引从5-9\%上调至9-12\%\par}
{\small\color{refgray}数据：资本开支从30-40亿升至50-70亿美元\par}
{\small\color{refgray}数据：开发回报约25\%\par}
{\small\color{refgray}数据：Q2新增9700个互联节点创纪录\par}
{\small\color{refgray}边际变化：此前市场关注AI对HBM的拉动，现发现AI需求已扩散至数据中心互联和传统存储，且Equinix以更高资本开支回应。\par}
\paragraph{MS} 中国生猪养殖产能出清慢于预期，生产效率提升（MSY提高）抵消能繁母猪下降，猪价反弹推迟。预计2026年均价11元/kg（同比-23\%），2027年回升至13.6元/kg，实质性反弹在2027年三季度。牧原成本优势仍存，但2026年预计全年亏损。
{\small\color{refgray}数据：2026年行业均价预计11元/kg（同比-23\%）\par}
{\small\color{refgray}数据：2027年均价预计13.6元/kg\par}
{\small\color{refgray}数据：猪价反弹推迟至2027年三季度\par}
{\small\color{refgray}数据：牧原完全成本11.5元/kg\par}
{\small\color{refgray}边际变化：此前市场预期2026年下半年周期性拐点，现因效率提升导致供给收紧延迟，拐点推迟至2027年。\par}
\paragraph{UBS} SK海力士自6月高点回调，但基本面未变。DRAM位元消耗量增速2027年升至36\%（2026年22\%），NAND升至23\%（20\%）。长期协议（LTA）已签10份，短期ASP承压但长期盈利稳定性改善。2027-28年EPS虽下调19\%，仍高于市场共识17\%以上。
{\small\color{refgray}数据：2027年DRAM位元消耗量增速36\%\par}
{\small\color{refgray}数据：2027年NAND位元消耗量增速23\%\par}
{\small\color{refgray}数据：LTA已签10份\par}
{\small\color{refgray}数据：2027-28年EPS下调19\%\par}
{\small\color{refgray}边际变化：此前市场关注存储周期见顶，现发现智能体AI拉动需求从HBM扩散至传统DRAM/NAND，且LTA改变盈利稳定性。\par}
\paragraph{Bernstein} 奔驰Q2调整后EBIT同比+16\%至22.99亿欧元，超预期40\%，但乘用车报告净利润仅4900万欧元（同比-94\%），因中国合资公司减值7.96亿欧元。欧洲和美国市场支撑业绩，中国以外BEV销量+51\%。管理层预计H2利润率弱于H1，全年在3-5\%区间下沿。
{\small\color{refgray}数据：调整后EBIT同比+16\%至22.99亿欧元\par}
{\small\color{refgray}数据：超卖方共识40\%\par}
{\small\color{refgray}数据：乘用车报告净利润4900万欧元（同比-94\%）\par}
{\small\color{refgray}数据：中国合资公司减值7.96亿欧元\par}
{\small\color{refgray}边际变化：此前市场关注经营韧性，现发现调整后利润与报告利润剪刀差巨大，中国业务减值对资产负债表影响显著。\par}
\paragraph{Bernstein} 美国消费者系统性转向价值导向，通胀3.5\%以上但失业率低位支撑消费韧性。沃尔玛、好市多、美元店受益；餐饮高端与平价两端分化，达登餐饮因消费者用餐厅牛排替代高价牛肉而受益。功能饮料增长较快，Celsius表现强劲。
{\small\color{refgray}数据：通胀率3.5\%以上\par}
{\small\color{refgray}数据：失业率低位\par}
{\small\color{refgray}数据：沃尔玛、好市多、美元店客流增长\par}
{\small\color{refgray}数据：达登餐饮受益于牛肉替代\par}
{\small\color{refgray}边际变化：此前市场关注消费整体韧性，现发现消费者支出方式结构性变化：价值导向渠道获得客流，高端与平价两端分化加剧。\par}
\paragraph{GS} Polaris Q2调整后EBITDA利润率同比回升540bp至11.8\%，剔除7400万美元关税退税后仍实现利润增长。北美ORV市场份额连续五个季度提升。对川崎重工具有参照意义：若获得类似退税，PS\&E利润存在向上修正空间。
{\small\color{refgray}数据：调整后EBITDA利润率同比+540bp至11.8\%\par}
{\small\color{refgray}数据：剔除7400万美元退税后仍增长\par}
{\small\color{refgray}数据：ORV零售单位销量同比+5\%\par}
{\small\color{refgray}数据：连续五个季度市场份额提升\par}
{\small\color{refgray}边际变化：此前市场关注退税一次性影响，现发现Polaris业务环境本身好转，北美四轮车市场定价和成本结构改善。\par}
\subsubsection*{关键数据}
\begin{itemize}
  \item 全球权益年初至今流入6590亿美元创纪录
  \item 中国Q2 GDP增速4.3\%，零售增速0.2\%
  \item 飞天茅台批价较底部涨13\%
  \item 2026年猪价预计11元/kg，同比-23\%
  \item 2027年DRAM位元消耗量增速36\%
  \item 奔驰中国合资公司减值7.96亿欧元
  \item 美国通胀3.5\%以上，失业率低位
  \item Polaris EBITDA利润率同比+540bp至11.8\%
\end{itemize}
\begin{figure}[htbp]
\centering
\includegraphics[width=0.88\linewidth]{figures/fig_001.jpg}
\caption{FIGURE 1 - \#\# Relative value - Underweight From a credit perspective, the balance sheet remains strong, with €30.4bn of net industrial liquidity after €5bn of shareholder distributions in H1. However, management expects H2 Cars profitability to be weaker than H1 due to h}
\end{figure}
\begin{figure}[htbp]
\centering
\includegraphics[width=0.88\linewidth]{figures/fig_007.jpg}
\caption{FIGURE 1 - \#\# Equity positioning near the highs \#\# Long only inflows were elevated again in July, albeit skewed by record China inflows FIGURE 2. Ytd equity inflows are now tracking ahead of record high bar set in 2021, at \textbackslash{}\$659bn FIGURE 3. Asset manager net futures posi}
\end{figure}
\begin{figure}[htbp]
\centering
\includegraphics[width=0.88\linewidth]{figures/fig_012.jpg}
\caption{EXHIBIT 2 - EXHIBIT 2: China GDP grew by 4.5\% YoY in Q2 EXHIBIT 3: The GDP deflator improved from -0.6\% in Q4 to -0.1\% in Q1 EXHIBIT 4: Export growth picked up in Q2}
\end{figure}
\begin{figure}[htbp]
\centering
\includegraphics[width=0.88\linewidth]{figures/fig_039.jpg}
\caption{Exhibit 3 - Exhibit 3: Muyuan – selling price vs. hog raising cost, by month Exhibit 4: Muyuan – hog sales volume, by month Exhibit 5: Hog prices in 2026: lowest point in 1H26, with rebound in 3Q26}
\end{figure}
{\small\color{refgray}本节综合 9 篇研究；涉及 BARC、Bernstein、MS、UBS、GS。}
\newpage
\subsection{AI与科技：从模型竞赛到生态博弈}
\textbf{AI产业正从模型参数竞赛转向生态与基础设施的深度博弈，企业采购重心从‘选模型’变为‘选生态’，边缘推理、先进封装与SaaS货币化成为下一阶段关键战场。}
\subsubsection*{主流共识}
\begin{itemize}
  \item AI基础设施需求强劲，先进封装设备订单连续两季环比增长超20\%，TCB与光通讯工具为核心驱动力。
  \item 企业AI采购从模型选择转向生态构建，微软Azure在产能受限下仍加速增长，验证AI单位经济改善。
  \item SaaS应用软件估值处于五年低位，AI货币化仍处早期，尚未成为收入加速的拐点。
  \item 边缘AI推理部署加速，Cloudflare等平台在分布式推理领域份额增长最快，成为新增长引擎。
\end{itemize}
\subsubsection*{分歧与条件差异}
\begin{itemize}
  \item AI模型规模定律是否持续：HSBC认为Kimi K3参数跃升至2.8万亿验证规模定律，但智谱GLM-5.2仅7440亿参数，追赶窗口收窄。
  \item AI货币化路径分歧：MS认为SaaS板块AI收入贡献仍为低个位数增量，而GS指出微软Copilot付费用户突破3000万，货币化路径日益清晰。
  \item ASEAN市场AI敞口分歧：UBS认为技术股仅占MSCI ASEAN 4\%权重却贡献25\%回报，估值分化下多数公司仍有空间；但外资流向分化，泰国马来西亚获净买入，印尼菲律宾持续流出。
\end{itemize}
\subsubsection*{机构观点}
\paragraph{GS} ASMPT订单连续两季环比增长45\%和24\%，3Q26预期高个位数环比增长，驱动力为TCB与光通讯工具，受益于数据中心CPO采用率上升、高端CPU/HBM需求及先进封装向面板级迁移。2026年营收预测上调至194亿港元，毛利率从40.3\%升至41.6\%，但净利润预测因税率影响基本不变。
{\small\color{refgray}数据：1Q26订单环比+45\%，2Q26环比+24\%\par}
{\small\color{refgray}数据：3Q26订单预期高个位数环比增长\par}
{\small\color{refgray}数据：2026年营收预测从181.51亿上调至194亿港元\par}
{\small\color{refgray}数据：毛利率从40.3\%上调至41.6\%\par}
{\small\color{refgray}边际变化：订单增长驱动力从传统封装转向TCB与光通讯工具，中国市场OSAT厂商向先进封装转型成为长期结构性增量。\par}
\paragraph{GS} 微软4QFY26营收900亿美元（+17\%），EPS 4.74美元（+23\%），Azure增速43\%超预期，在GPU产能约束下仍加速增长。Copilot付费用户突破3000万，单季净增1000万。企业AI采购重心从‘选模型’转向‘选生态’，微软角色更趋战略性，GPU上线时间缩短近50\%，定价改进与自研芯片有望推动FY27利润率阶跃。
{\small\color{refgray}数据：4QFY26营收900亿美元（+17\%），EPS 4.74美元（+23\%）\par}
{\small\color{refgray}数据：Azure恒定汇率增速43\%，超市场预期40\%\par}
{\small\color{refgray}数据：Copilot付费用户突破3000万，单季净增1000万\par}
{\small\color{refgray}数据：GPU上线时间缩短近50\%\par}
{\small\color{refgray}边际变化：微软从模型提供商转向生态中枢，AI单位经济改善且云利润率稳定，市场对AI投入回报的担忧得到缓解。\par}
\paragraph{UBS} MSCI ASEAN技术股仅占4\%权重却贡献25\%回报，55\%成份股估值低于10年均值，指数昂贵掩盖广泛公司吸引力。外资流向分化：泰国和马来西亚因改革动能获净买入，印尼和菲律宾持续流出。AI直接敞口有限，但数据中心、电力等供应链环节仍有机会，利率‘更高更久’环境下菲律宾对美利率最敏感。
{\small\color{refgray}数据：技术股占MSCI ASEAN权重4\%，贡献25\%回报\par}
{\small\color{refgray}数据：55\%成份股估值低于10年均值\par}
{\small\color{refgray}数据：ASEAN自KOSPI 2026年6月见顶以来跑赢新兴市场16\%\par}
{\small\color{refgray}数据：泰国和马来西亚获外资净买入，印尼菲律宾持续流出\par}
{\small\color{refgray}边际变化：市场对ASEAN的‘AI敞口不足’叙事被估值分化数据挑战，分散化需求可能成为重新关注催化剂。\par}
\paragraph{MS} 应用软件2Q26需求基本面稳固，但AI货币化仍处早期，管理层维持谨慎展望，广泛上调修正推动板块重估可能性有限。SaaS应用组EV/S倍数低于五年均值，相对软件板块折价45\%接近历史低位。Figma首个完整AI信用货币化季度超75\%企业用户继续消费，HubSpot AI采用率约20-25\%，信用额度贡献低个位数收入。
{\small\color{refgray}数据：SaaS应用组EV/S倍数低于五年均值\par}
{\small\color{refgray}数据：相对软件板块折价45\%\par}
{\small\color{refgray}数据：Figma超75\%组织企业用户在免费额度用完后继续消费\par}
{\small\color{refgray}数据：HubSpot AI采用率约20-25\%，信用额度贡献低个位数\par}
{\small\color{refgray}边际变化：AI货币化从概念验证进入早期收入贡献阶段，但尚未成为收入加速的拐点，低估值本身不构成买入理由。\par}
\paragraph{MS} Cloudflare正从CDN安全公司演变为边缘AI推理基础设施关键参与者。渠道调研显示其仍是合作伙伴增长最快供应商之一，边缘AI调查中55\%合作伙伴认为其过去三个月边缘推理市场份额增长最快，63\%选为边缘AI计算首选平台。Q2营收预计约6.65亿美元（+29.8\%），cRPO同比增长预计达33\%。
{\small\color{refgray}数据：Q2营收预计约6.65亿美元（同比+29.8\%）\par}
{\small\color{refgray}数据：cRPO同比增长预计达33\%\par}
{\small\color{refgray}数据：55\%合作伙伴认为Cloudflare边缘推理市场份额增长最快\par}
{\small\color{refgray}数据：63\%选为边缘AI计算首选平台\par}
{\small\color{refgray}边际变化：市场认知从CDN安全转向边缘AI推理，企业级大规模部署预计2027年，但早期推理工作负载已开始贡献收入。\par}
\paragraph{MS} 金山办公1H26收入预计同比增长21-28\%，2Q26单季增长18-33\%，对应16-18亿人民币，增长来自WPS个人订阅和WPS 365企业协作产品。调整后利润同比增长3-45\%，中位数约24\%与收入增速一致，利润率保持稳定。当前2026/2027年PE分别为63/46倍，估值不便宜。
{\small\color{refgray}数据：1H26收入预计同比增长21-28\%\par}
{\small\color{refgray}数据：2Q26单季收入16-18亿人民币\par}
{\small\color{refgray}数据：调整后利润中位数同比增长约24\%\par}
{\small\color{refgray}数据：2026/2027年PE分别为63/46倍\par}
{\small\color{refgray}边际变化：收入增长缓解市场对AI冲击传统办公软件的担忧，但估值偏高是主要顾虑，需关注增长来自存量付费转化还是新场景拓展。\par}
\paragraph{HSBC} 智谱AI 2026年7月ARR达10亿美元，比指引提前约5个月，HSBC上调年底ARR至20亿美元，2027年底至65亿美元。但收入结构向低毛利API倾斜，竞争加剧致定价承压，H股增发后股本扩张。Kimi K3参数从1万亿跃升至2.8万亿，验证规模定律，智谱GLM-5.2仅7440亿参数，追赶窗口收窄。预计2027年盈利。
{\small\color{refgray}数据：2026年7月ARR达10亿美元，提前约5个月\par}
{\small\color{refgray}数据：2026年底ARR预测从10亿上调至20亿美元\par}
{\small\color{refgray}数据：Kimi K3总参数2.8万亿，激活参数1040亿\par}
{\small\color{refgray}数据：智谱GLM-5.2总参数7440亿\par}
{\small\color{refgray}边际变化：ARR提前达标但竞争加剧，模型参数差距扩大，盈利时间提前但FCF假设下调导致目标价从1900降至1500港元。\par}
\paragraph{Bernstein} Air Liquide 1H26业绩为‘减速带’而非趋势转折，订单积压创纪录60亿欧元，环比新增5亿，未来2-3年收入可见度高。电子业务预计保持8-9\%可持续增长，2H26增速上调40个基点。氦气供应短期影响可控，卡塔尔生产已恢复但出口受限需几个季度恢复正常。维持Outperform，目标价从189上调至202欧元。
{\small\color{refgray}数据：订单积压创纪录60亿欧元，环比新增5亿\par}
{\small\color{refgray}数据：电子业务预计8-9\%可持续增长\par}
{\small\color{refgray}数据：2H26电子业务增长预期上调40个基点\par}
{\small\color{refgray}数据：目标价从189上调至202欧元\par}
{\small\color{refgray}边际变化：短期氦气扰动被市场过度解读，订单积压新高和电子业务加速为下半年增长提供支撑，10月资本市场日为关键催化剂。\par}
\subsubsection*{关键数据}
\begin{itemize}
  \item ASMPT 1Q26订单环比+45\%，2Q26环比+24\%
  \item 微软Azure恒定汇率增速43\%，Copilot付费用户突破3000万
  \item MSCI ASEAN技术股仅4\%权重贡献25\%回报，55\%成份股估值低于10年均值
  \item SaaS应用组EV/S倍数低于五年均值，相对软件板块折价45\%
  \item Cloudflare Q2营收预计6.65亿美元（+29.8\%），63\%合作伙伴选为边缘AI首选平台
  \item 智谱AI ARR提前5个月达10亿美元，Kimi K3参数跃升至2.8万亿
  \item Air Liquide订单积压创纪录60亿欧元，电子业务预计8-9\%增长
\end{itemize}
\begin{figure}[htbp]
\centering
\includegraphics[width=0.88\linewidth]{figures/fig_079.jpg}
\caption{Exhibit 1 - Exhibit 1: ASMPT's BB ratio was 1.43 by 2Q26 (vs. 1.43 by 1Q26) Exhibit 2: ASMPT's revenues: Recovering trend Earnings revision: We factor in ASMPT's 2Q26 result and raise our 2026E revenue and GM to reflect the stronger the ex}
\end{figure}
\begin{figure}[htbp]
\centering
\includegraphics[width=0.88\linewidth]{figures/fig_111.jpg}
\caption{Exhibit 1 - Exhibit 1: Comparison between Z.ai and Moonshot on key metrics}
\end{figure}
\begin{figure}[htbp]
\centering
\includegraphics[width=0.88\linewidth]{figures/fig_126.jpg}
\caption{Exhibit 1 - Exhibit 1: Overall Software Multiples and SaaS Applications Group Both Trading Below 5-Year Trailing Average Exhibit 2: SaaS Applications Software Group EV/S Multiples Trading at a 45\% Discount to Broader Software Exhibit 3: EV}
\end{figure}
\begin{figure}[htbp]
\centering
\includegraphics[width=0.88\linewidth]{figures/fig_131.jpg}
\caption{Exhibit 9 - Exhibit 3: NET: MSe vs. Consensus}
\end{figure}
{\small\color{refgray}本节综合 8 篇研究；涉及 Bernstein、UBS、GS、HSBC、MS。}
\newpage
\subsection{AI驱动结构性分化，存储与被动元件定价权增强}
\textbf{半导体与硬件板块呈现显著的结构性分化：AI推理芯片测试需求、HDD定价韧性和被动元件产能趋紧构成上行主线，而汽车模拟芯片和传统设备商则面临增长节奏调整。存储定价在供应约束下保持强势，HBM弹性可能被低估。}
\subsubsection*{主流共识}
\begin{itemize}
  \item AI推理芯片测试需求成为设备市场新增长引擎，爱德万测试上调SoC测试设备市场预期19亿美元。
  \item HDD需求可见性从2027年延伸至2028年，客户范围从云服务商扩展至AI模型开发者与机器人领域。
  \item 被动元件产能利用率逼近满载，AI营收占比提前突破全年目标，订单出货比升至2.2。
  \item 存储DRAM定价全年保持双位数季度环比增长，长期协议约束力增强，价格波动性结构性下降。
\end{itemize}
\subsubsection*{分歧与条件差异}
\begin{itemize}
  \item HBM定价弹性分歧：GS预计三星HBM定价明年同比增长87\%，高于彭博一致预期的52\%，源于常规DRAM传导效应。
  \item 汽车半导体复苏节奏分歧：Bernstein认为NXP汽车业务面临中国调整信号，而HSBC认为意法半导体复苏仅是节奏调整，趋势未变。
  \item 设备商短期业绩分化：KLA 2026年因产品组合成为份额让渡方，而爱德万测试和Seagate持续超预期。
\end{itemize}
\subsubsection*{机构观点}
\paragraph{Bernstein} NXP Q2营收34.96亿美元、EPS 3.61美元均小幅超预期，通信基础设施和工业物联网是主要超预期来源，渠道库存稳定在11周。但汽车终端市场尤其是中国出现调整信号，数据中心业务占比仅约3\%，远低于同行。估值在板块回调中变得有吸引力，但上行空间取决于市场如何定价这些结构性矛盾。
{\small\color{refgray}数据：Q2营收34.96亿美元（预期34.63亿）\par}
{\small\color{refgray}数据：EPS 3.61美元（预期3.54）\par}
{\small\color{refgray}数据：毛利率58.0\%符合预期\par}
{\small\color{refgray}数据：渠道库存11周，处于长期目标水平\par}
{\small\color{refgray}边际变化：Q3指引略超预期，但汽车业务增长叙事面临中国需求调整压力，AI故事缺失成为估值折价核心。\par}
\paragraph{MS} Seagate F4Q26财报显示定价是驱动收入、利润率、EPS和FCF超预期的主要因素，且趋势尚未见顶。需求可见性从贯穿2027年延伸至2028年，客户范围从云服务商扩展至AI模型开发者、机器人等。9月季度定价同比+25\%，毛利率从47\%提升至52.7\%，增量毛利率83\%。CY27 EPS预期从35.75美元上调至59.33美元。
{\small\color{refgray}数据：9月季度定价同比+25\%（6月季度+10\%）\par}
{\small\color{refgray}数据：毛利率从47\%提升至52.7\%，下季度指引约57.5\%\par}
{\small\color{refgray}数据：增量毛利率83\%（上季度72\%）\par}
{\small\color{refgray}数据：CY27 EPS预期从35.75美元上调至59.33美元\par}
{\small\color{refgray}边际变化：需求可见性从2027年延伸至2028年，客户类型从传统云向AI模型开发者、机器人等扩展，HDD需求基础正在多元化。\par}
\paragraph{MS} KLA 2026财年业绩呈现混合信号，9月季度指引略低于预期，未能实现beat and raise。2026年产品组合对工艺控制强度形成影响，KLA成为份额让渡方。但2027年结构性故事更清晰：125亿美元RPO（FY25年底78.6亿），1D/1.4nm节点和先进封装驱动工艺控制强度回升，预计2027年营收增长36\%，跑赢WFE的31\%。
{\small\color{refgray}数据：2026年WFE预测从1400亿美元上调至1500亿美元以上\par}
{\small\color{refgray}数据：2026年先进封装收入预期从10亿上调至11亿美元（同比+70\%）\par}
{\small\color{refgray}数据：RPO 125亿美元（FY25年底78.6亿）\par}
{\small\color{refgray}数据：预计2027年营收增长36\%，跑赢WFE的31\%\par}
{\small\color{refgray}边际变化：2026年短期业绩受产品组合拖累，但2027年工艺控制强度回升逻辑更清晰，RPO大幅增长提供远期可见性。\par}
\paragraph{GS} 爱德万测试Q1营业利润1900亿日元，超出GS预期约25\%，全年指引从6275亿日元大幅上调至8460亿日元。AI推理芯片测试需求成为新增长引擎，SoC测试设备市场总规模预期上调19亿美元。公司SoC测试设备市场份额从66\%持续上升，中期目标至少70\%。毛利率创历史新高69.5\%。
{\small\color{refgray}数据：Q1营业利润1900亿日元（超GS预期25\%）\par}
{\small\color{refgray}数据：全年营业利润指引从6275亿上调至8460亿日元\par}
{\small\color{refgray}数据：SoC测试设备市场总规模预期上调19亿美元\par}
{\small\color{refgray}数据：毛利率创历史新高69.5\%\par}
{\small\color{refgray}边际变化：AI推理芯片测试需求成为新增长引擎，而非仅依赖GPU测试时间延长，市场对业绩超预期准备不足。\par}
\paragraph{GS} MARUWA Q1营业利润63亿日元，低于GS预期70亿日元，但全年指引从297亿上调至337亿日元。Q1利润缺口源于低利润率AI相关库存出货，已在Q1基本消化完毕，良率问题也已解决。AI应用全年销售指引弹性较高，通信领域AI收入约60亿日元，半导体应用需求恢复提前至Q2，订单创历史新高。
{\small\color{refgray}数据：Q1营业利润63亿日元（GS预期70亿）\par}
{\small\color{refgray}数据：全年营业利润指引从297亿上调至337亿日元\par}
{\small\color{refgray}数据：通信领域Q1收入约92亿日元，其中AI约60亿\par}
{\small\color{refgray}数据：全年AI收入指引约310亿日元\par}
{\small\color{refgray}边际变化：Q1利润低于预期但全年指引上调，低利润率库存出清后出货结构将转向更高利润率产品组合，AI订单确认提供可见性。\par}
\paragraph{GS} 韩国存储专家会议判断DRAM定价全年保持双位数季度环比增长，4季度仍有再次双位数增长可能，核心驱动力是供应短缺。HBM定价明年弹性可能被低估，GS预计三星HBM定价同比增长87\%，高于彭博一致预期的52\%。超过一半的服务器DRAM已通过长期协议覆盖，包含大额预付款和照付不议条款。
{\small\color{refgray}数据：DRAM 3季度定价双位数季度环比增长\par}
{\small\color{refgray}数据：4季度仍有双位数季度环比增长可能\par}
{\small\color{refgray}数据：GS预计三星HBM定价明年同比+87\%（彭博一致预期52\%）\par}
{\small\color{refgray}数据：超过一半的服务器DRAM已签署LTA\par}
{\small\color{refgray}边际变化：HBM定价上行空间可能被市场低估，常规DRAM定价持续走强对HBM定价有传导效应，LTA覆盖率提升降低价格波动性。\par}
\paragraph{GS} 国巨Q2 AI营收占比达16\%，提前超越全年15\%目标。产能利用率80-85\%，预计Q3提升至90\%以上，整体产能接近满载。订单出货比逐月改善至6月底的2.2（Q1仅1.0+），是下半年需求走强的明确信号。客户长约需求增加，尤其集中在MLCC和钽电容。但基于MLCC和钽电容定价假设调整，下调2026-2028年盈利预测10/17/14\%。
{\small\color{refgray}数据：AI营收占比16\%（Q1为14-15\%），提前超越全年15\%目标\par}
{\small\color{refgray}数据：产能利用率80-85\%，Q3预计提升至90\%以上\par}
{\small\color{refgray}数据：订单出货比6月底达2.2（Q1仅1.0+）\par}
{\small\color{refgray}数据：下调2026/27/28年盈利预测10/17/14\%\par}
{\small\color{refgray}边际变化：AI营收占比提前达标，订单出货比大幅攀升显示需求走强，但MLCC和钽电容定价假设调整导致盈利预测下调。\par}
\paragraph{HSBC} 意法半导体Q2订单出货比接近2，所有终端市场均高于1，通信设备和电脑外设板块显著高于2，分销库存已低于公司目标水平。但Q3营收指引低于预期，复苏节奏慢于预期。2026年AI数据中心收入预期上调至超过10亿美元，2027年远超20亿美元。下调2026/2027年营业利润预期4.4\%和5.9\%，维持2028/2030年EPS 4.1美元和6.2美元预测。
{\small\color{refgray}数据：整体订单出货比接近2\par}
{\small\color{refgray}数据：通信设备板块订单出货比显著高于2\par}
{\small\color{refgray}数据：分销库存已低于公司目标水平\par}
{\small\color{refgray}数据：2026年AI数据中心收入预期超过10亿美元\par}
{\small\color{refgray}边际变化：Q3指引低于预期但订单数据强劲，复苏节奏调整而非趋势中断，AI/光互联成为最大增量。\par}
\subsubsection*{关键数据}
\begin{itemize}
  \item 爱德万测试全年营业利润指引从6275亿上调至8460亿日元
  \item Seagate 9月季度定价同比+25\%，增量毛利率83\%
  \item 国巨订单出货比6月底达2.2（Q1仅1.0+）
  \item GS预计三星HBM定价明年同比+87\%（彭博一致预期52\%）
  \item KLA RPO 125亿美元（FY25年底78.6亿）
  \item 意法半导体订单出货比接近2，分销库存低于目标
  \item NXP数据中心业务占比仅约3\%
  \item MARUWA全年AI收入指引约310亿日元
\end{itemize}
\begin{figure}[htbp]
\centering
\includegraphics[width=0.88\linewidth]{figures/fig_047.jpg}
\caption{Exhibit 1 - Exhibit 1: STX F4Q26 Earnings Variance}
\end{figure}
\begin{figure}[htbp]
\centering
\includegraphics[width=0.88\linewidth]{figures/fig_048.jpg}
\caption{Exhibit 1 - Exhibit 1: STX F4Q26 Earnings Variance SEAGATE \#\# Financial Model}
\end{figure}
\begin{figure}[htbp]
\centering
\includegraphics[width=0.88\linewidth]{figures/fig_100.jpg}
\caption{Exhibit 1 - Exhibit 2: Yageo revenue mix by channel (2Q26) Exhibit 3: Yageo revenue mix by region (2Q26) Exhibit 4: Yageo revenue mix by product (2Q26)}
\end{figure}
\begin{figure}[htbp]
\centering
\includegraphics[width=0.88\linewidth]{figures/fig_136.jpg}
\caption{Exhibit 1 - \#\# OWNERSHIP POSITIONING Refinitiv; MSPB Content. Includes certain hedge fund exposures held with MSPB. Information may be inconsistent with or may not reflect broader market trends. Long/Short Ratio = Long Exposure / Short exposure. Sector \% of Total Net Expo}
\end{figure}
{\small\color{refgray}本节综合 8 篇研究；涉及 Bernstein、MS、GS、HSBC。}
\newpage
\subsection{消费与零售：分化加剧，品牌力与渠道重塑成为核心主线}
\textbf{全球消费与零售市场呈现显著分化：高端奢侈品受益于高净值客群韧性，但大众消费与入门级需求承压；运动品牌与餐饮企业面临增长引擎切换与利润稀释的挑战；中国婴幼儿配方奶粉行业受质量事件冲击，但头部品牌加速整合。品牌力、渠道效率与品类组合成为决定企业表现的关键变量。}
\subsubsection*{主流共识}
\begin{itemize}
  \item 高端奢侈品市场韧性主要来自高净值客群，大众消费与入门级需求持续疲软。
  \item 中国消费市场结构性增长空间仍在，但动力从大众消费者转向高净值人群，且消费行为更趋挑剔、注重体验与本地偏好。
  \item 运动品牌正从潮流驱动转向功能驱动，但利润端受营销与履约成本上升挤压。
\end{itemize}
\subsubsection*{分歧与条件差异}
\begin{itemize}
  \item 爱马仕的品类组合向上延伸（高级珠宝、高级定制）能否有效对冲皮革制品产量增长放缓，市场存在分歧。
  \item Nike中国线上渠道调整（终止批发合作）的短期收入损失与长期品牌价值提升之间的权衡，机构看法不一。
  \item 海底捞核心业务恢复节奏慢于预期，新业务（外卖、子品牌）对利润率的稀释程度，市场尚未形成一致预期。
  \item 中国婴幼儿配方奶粉行业质量事件的影响范围与持续时间，以及头部品牌能否持续扩大份额，存在不确定性。
\end{itemize}
\subsubsection*{机构观点}
\paragraph{Bernstein} 爱马仕1H26有机销售+6.7\%符合预期，但EBIT利润率41.1\%超预期70bps，毛利率提升是主因。区域分化明显：法国（+6.2\%）与日本（+12.3\%）超预期，亚太（除日本）仅+2.5\%略低于预期。品类上，丝绸（+12.2\%）与手表（+4.4\%）意外走强，香水（-9.5\%）拖累。市场过度聚焦皮革制品产量增长（+6\%为生产工时而非产量），忽略了品类组合向上延伸（高级珠宝、高级定制）与门店扩张带来的增长空间。账上现金超120亿欧元，特别分红概率上升。
{\small\color{refgray}数据：1H26有机销售+6.7\%，符合预期\par}
{\small\color{refgray}数据：EBIT利润率41.1\%，超预期70bps\par}
{\small\color{refgray}数据：法国Q2有机增长+6.2\%，预期+3.3\%\par}
{\small\color{refgray}数据：日本Q2有机增长+12.3\%，预期+10.6\%\par}
{\small\color{refgray}边际变化：市场焦点从皮革制品产量增长转向品类组合与门店扩张的增量贡献，利润率超预期成为新亮点。\par}
\paragraph{NOM} 爱马仕2Q26亚太（除日本）收入同比+2.5\%，为所有区域最慢，低于预期+4.5\%。但管理层确认中国区销售仍在增长且出现企稳迹象。高净值客户韧性较强，而“有抱负的消费者”面临更大挑战，高价位产品（珠宝、成衣）表现优于香水（-9.5\%）。开云集团2Q26可比收入同比+2\%，略超预期，核心品牌Gucci降幅从-8\%收窄至-2\%。亚太地区收入降幅从-4\%收窄至-1\%，但中国大陆仍在调整，韩国表现突出。中国消费者正变得更挑剔、更注重体验、更倾向本土品牌。
{\small\color{refgray}数据：爱马仕2Q26亚太（除日本）收入同比+2.5\%，预期+4.5\%\par}
{\small\color{refgray}数据：爱马仕香水美容同比-9.5\%\par}
{\small\color{refgray}数据：开云2Q26可比收入同比+2\%，预期36.2亿欧元\par}
{\small\color{refgray}数据：Gucci 2Q26可比收入同比-2\%，1Q26为-8\%\par}
{\small\color{refgray}边际变化：中国消费者行为结构性转变（更挑剔、更重体验、更倾向本土品牌）成为国际品牌调整策略的关键驱动。\par}
\paragraph{Bernstein} 阿迪达斯2Q26营收货币中性增长14\%，超预期13\%，但营业利润率8.5\%低于预期9.4\%，主因营销费用（世界杯、产品发布）与DTC履约成本超支。运动表现品类增长39\%（1Q26为29\%），而生活方式仅增2\%。鞋类仅增1\%，服装增35\%。北美（+17\%）与拉美（+28\%）增速加快，欧洲批发仍保守。Nike中国线上渠道调整：2027年1月起终止合作伙伴线上店铺，仅保留官方直营与旗舰店，预计带来约10亿美元收入下滑，但毛利率有望改善约200bps。Nike中国份额已从27\%降至16\%。
{\small\color{refgray}数据：阿迪达斯2Q26营收货币中性+14\%，预期+13\%\par}
{\small\color{refgray}数据：阿迪达斯2Q26营业利润率8.5\%，预期9.4\%\par}
{\small\color{refgray}数据：阿迪达斯运动表现品类+39\%，生活方式+2\%\par}
{\small\color{refgray}数据：阿迪达斯鞋类+1\%，服装+35\%\par}
{\small\color{refgray}边际变化：阿迪达斯增长引擎从潮流转向功能，但利润端受成本上升挤压；Nike中国渠道调整从短期增长转向品牌质量与定价能力。\par}
\paragraph{GS} 中国婴幼儿配方奶粉行业7月连发两起质量事件：Cabio因ARA问题发布盈利预警，预计1H26净亏损约1.01亿元，触发ST风险警示；Friso澳门批次铅含量超标（第三方复检<0.007mg/kg，远低于澳门限值0.02mg/kg），但百度搜索热度峰值远低于年初雀巢事件。行业线上销售连续两季度下滑，但头部品牌如贝拉米（+70\%）与Biostime（+43\%）6月线上销售额同比大增，显示强品质管控与品牌力企业正在扩大份额。
{\small\color{refgray}数据：Cabio预计1H26净亏损约1.01亿元\par}
{\small\color{refgray}数据：Friso澳门批次铅含量超标，第三方复检<0.007mg/kg\par}
{\small\color{refgray}数据：Friso事件百度搜索热度峰值远低于年初雀巢事件\par}
{\small\color{refgray}数据：贝拉米6月线上销售额同比+70\%\par}
{\small\color{refgray}边际变化：行业从需求调整与监管趋严转向质量事件冲击下的格局分化，头部品牌加速整合。\par}
\paragraph{GS} 古茗控股预计1H26营收同比+30\%，调整后核心净利润同比+43\%，但2026年净新增门店预期从3600家下调至2000家，主因加盟商意愿减弱与公司优先门店质量。全年单店GMV预计同比-2\%。自2Q26高点股价回调19\%，当前交易在13倍2026年预期PE。海底捞1H26收入预计同比+8\%，但净利润仅+5\%，主因翻台率增长仅约2\%、客单价下滑约1\%、门店关闭与新业务稀释利润率。全年单店收入预期从+3\%下调至-0.5\%。
{\small\color{refgray}数据：古茗1H26预计营收同比+30\%\par}
{\small\color{refgray}数据：古茗1H26预计调整后核心净利润同比+43\%\par}
{\small\color{refgray}数据：古茗2026年净新增门店预期从3600家降至2000家\par}
{\small\color{refgray}数据：古茗当前交易在13倍2026年预期PE\par}
{\small\color{refgray}边际变化：古茗从数量扩张转向质量提升，海底捞核心业务从量价齐升转向量稳价降，新业务尚不能弥补缺口。\par}
\paragraph{JPM} 截至7月29日当周，零售资金流降至过去一年历史百分位的0.8\%，为年内最弱水平之一。散户连续四天净卖出个股，仅小幅买入ETF。科技硬件股遭遇有记录以来最大单周资金流出，半导体板块拥挤度从99百分位降至91.8百分位。科技ETF流出规模为历史第二，主要驱动力来自三倍做多半导体ETF（SOXL）。美光（MU）从散户最青睐名单跌入最卖名单。
{\small\color{refgray}数据：零售资金流降至过去一年历史百分位的0.8\%\par}
{\small\color{refgray}数据：科技硬件股单周卖出规模创纪录\par}
{\small\color{refgray}数据：半导体板块拥挤度从99百分位降至91.8百分位\par}
{\small\color{refgray}数据：科技ETF流出规模为历史第二\par}
{\small\color{refgray}边际变化：散户从逢跌买入转为连续净卖出，科技股动量反转与去拥挤化成为主导，半导体板块调整尤为显著。\par}
\subsubsection*{关键数据}
\begin{itemize}
  \item 爱马仕1H26 EBIT利润率41.1\%，超预期70bps
  \item 爱马仕亚太（除日本）Q2有机增长+2.5\%，低于预期+4.5\%
  \item Gucci 2Q26可比收入同比-2\%，1Q26为-8\%
  \item 阿迪达斯2Q26营业利润率8.5\%，低于预期9.4\%
  \item Nike中国份额从27\%降至16\%
  \item Cabio预计1H26净亏损约1.01亿元
  \item 古茗2026年净新增门店预期从3600家降至2000家
  \item 海底捞全年单店收入预期从+3\%下调至-0.5\%
  \item 科技硬件股单周卖出规模创纪录
  \item 半导体板块拥挤度从99百分位降至91.8百分位
\end{itemize}
\begin{figure}[htbp]
\centering
\includegraphics[width=0.88\linewidth]{figures/fig_022.jpg}
\caption{EXHIBIT 3 - EXHIBIT 3: Luxury Goods Calendar Year 2Q26 results CY 2Q26 Reported Organic Growth and Actual Change in Revenue \#\# INVESTMENT IMPLICATIONS \#\# We rate Hermès Outperform with a target price of EUR2,150.}
\end{figure}
\begin{figure}[htbp]
\centering
\includegraphics[width=0.88\linewidth]{figures/fig_064.jpg}
\caption{EXHIBIT 1 - EXHIBIT 1: Nike's share has declined from around 27\% at its peak to around 16\% in 2025, with much of the pre-COVID share expansion now having been reversed EXHIBIT 2: Since 2020, local competitors have steadily taken share from I}
\end{figure}
\begin{figure}[htbp]
\centering
\includegraphics[width=0.88\linewidth]{figures/fig_087.jpg}
\caption{Exhibit 1 - Exhibit 1: MNC players saw sequential market share trend mixed in Mar-Apr, Friesland/Danone/A2 value share increased while Wyeth decreased vs. Jan-Feb Offline international IMF brands market share by value Exhibit 2: IMF 618 Shop}
\end{figure}
\begin{figure}[htbp]
\centering
\includegraphics[width=0.88\linewidth]{figures/fig_116.jpg}
\caption{Figure 54 - Figure 1: Momentum Crowding Figure 2: Momentum Unwind Now vs Feb-Mar'25 Figure 3: Crowding in Low Quality}
\end{figure}
{\small\color{refgray}本节综合 10 篇研究；涉及 Bernstein、NOM、GS、JPM。}
\newpage
\subsection{工业与能源：电网投资与AI驱动下的结构性分化}
\textbf{全球工业与能源板块正经历结构性分化：电网投资（变压器、能源管理）和AI基础设施（高端CCL、燃料电池）需求强劲，而传统化工、汽车等周期行业仍面临量强价弱、现金流承压的挑战。中国供应商在变压器和继电器领域的全球份额持续提升，但美国政策风险（如FCC逆变器限制）为部分企业带来不确定性。}
\subsubsection*{主流共识}
\begin{itemize}
  \item 全球电网投资持续高景气，变压器供需缺口（美国当前72\%）预计到2028年才收窄至57\%，中国供应商凭借短交货周期（24-36周 vs 全球128周）填补缺口。
  \item AI数据中心需求是工业增长的核心驱动力，施耐德电气数据中心业务实现三位数增长，高端CCL（M8/M9级）产能已被提前锁定至2028年。
  \item 工业自动化在中国出现复苏信号，施耐德中国区增速19.6\%，与西门子正向共振，但持续性仍需观察订单转化率。
\end{itemize}
\subsubsection*{分歧与条件差异}
\begin{itemize}
  \item 美国对中国逆变器限制（FCC DA 26-786）采用《购买美国产品法》测试而非原产地规则，第三国产能无法规避，这与市场此前预期存在分歧。
  \item 化工行业（巴斯夫）量强价弱，销量超预期但价格低于预期，现金流大幅恶化（现金转换率从89\%降至21\%），而继电器（宏发）和变压器（中国出口）则通过提价成功传导成本。
  \item 力拓的铜当量产量增长路径和成本削减力度（目标18亿美元）比市场预期更激进，但项目评估支出增加（同比+122\%）显示增长与成本控制之间的平衡。
\end{itemize}
\subsubsection*{机构观点}
\paragraph{Bernstein} 霍尼韦尔二季度有机收入同比+4\%（预期0\%），调整后EPS 1.95美元（超共识7\%），三大业务板块均超预期。订单增长16\%，订单出货比>1.1倍，下半年增长指引4-6\%偏保守。工业自动化利润率从17\%到22\%的路径清晰，一半来自剥离低毛利业务，一半来自经营杠杆。市场对管理层可信度的评估正在重新调整。
{\small\color{refgray}数据：有机收入同比+4\%，市场预期0\%\par}
{\small\color{refgray}数据：调整后EPS 1.95美元，超共识7\%\par}
{\small\color{refgray}数据：总订单有机增长16\%，各板块两位数增长\par}
{\small\color{refgray}数据：工业自动化利润率目标22\%，路径清晰\par}
{\small\color{refgray}边际变化：此前持怀疑态度的分析师开始重新评估，二季度数据使'展示给我看'故事变得可信。\par}
\paragraph{NOM} FCC将外国产逆变器纳入覆盖清单，阳光电源产品组合大概率被覆盖，储能PCS也受影响。标准采用《购买美国产品法》测试，第三国产能同样受限。现有产品通过祖父条款可延续至2026年，但新机型认证通道关闭。长期看，固态变压器等第二增长曲线叙事被削弱，合规成本上升。下调2027/28年EPS至8.78/10.41元。
{\small\color{refgray}数据：FCC DA 26-786将逆变器纳入覆盖清单\par}
{\small\color{refgray}数据：采用Buy America测试，非原产地规则\par}
{\small\color{refgray}数据：2027/28年EPS下调至8.78/10.41元（原9.13/11.20元）\par}
{\small\color{refgray}数据：目标价120元，基于14倍2027年PE\par}
{\small\color{refgray}边际变化：市场此前认为东南亚/印度建厂可规避限制，但FCC援引Buy America测试，第三国产能策略合规性存根本疑问。\par}
\paragraph{GS} 中国变压器出口结构性分化：6月总出口同比+47\%（5月+72\%），但对美出口同比+172\%、环比+52\%。美国供需缺口当前72\%，预计2028年收窄至57\%，全球交货周期128周，中国供应商仅24-36周。中东（+141\%）和美洲（+108\%）是主要增量来源，亚洲转弱（-38\%）。220-330MVA段出口均价近3月滚动同比-8\%，但6月单月已转正为+3\%。
{\small\color{refgray}数据：6月中国变压器出口同比+47\%，对美出口+172\%\par}
{\small\color{refgray}数据：美国供需缺口72\%，预计2028年收窄至57\%\par}
{\small\color{refgray}数据：全球交货周期128周，中国供应商24-36周\par}
{\small\color{refgray}数据：中东出口+141\%，亚洲出口-38\%\par}
{\small\color{refgray}边际变化：对美出口增速与整体增速背离，美国供需紧张未缓解，中国供应商仍在填补缺口。\par}
\paragraph{GS} 巴斯夫二季度调整后EBITDA 24.49亿欧元，超预期18\%，但量强价弱：销量同比+7.3\%（预期+2.0\%），价格仅+11.5\%（预期+12.7\%）。运营现金流仅5.24亿欧元（去年同期15.85亿），现金转换率从89\%骤降至21\%，主因转型支出和营运资本增加。材料和工业解决方案是利润超预期主要来源，但化学品板块利润率低于预期486bps。
{\small\color{refgray}数据：调整后EBITDA 24.49亿欧元，超预期18\%\par}
{\small\color{refgray}数据：销量同比+7.3\%，远超预期+2.0\%\par}
{\small\color{refgray}数据：运营现金流5.24亿欧元，现金转换率21\%\par}
{\small\color{refgray}数据：化学品板块利润率低于预期486bps\par}
{\small\color{refgray}边际变化：现金流数据揭示了转型真实成本，市场此前对盈利质量关注不足。\par}
\paragraph{GS} 高端CCL供需从'产能紧张'转向'产能被提前锁定'，毛利率扩张驱动力切换为定价权提升与产品结构改善。台光2Q26毛利率环比+4.4ppt至33.9\%，台燿+4.6ppt至29.7\%，均超彭博一致预期。3Q26展望：Trainium 3/Rubin等新AI项目放量，台光毛利率预计升至36.8\%，台燿再扩3.4ppt。2028年后产能已被提前预订，两家公司未来两年扩产35-50\%+。
{\small\color{refgray}数据：台光2Q26毛利率33.9\%，环比+4.4ppt\par}
{\small\color{refgray}数据：台燿2Q26毛利率29.7\%，环比+4.6ppt\par}
{\small\color{refgray}数据：3Q26台光毛利率预计升至36.8\%\par}
{\small\color{refgray}数据：未来两年扩产35-50\%+，产能已被提前锁定\par}
{\small\color{refgray}边际变化：毛利率扩张从一次性事件变为结构性变化，基板级CCL认证通过后毛利率可达40-50\%+。\par}
\paragraph{GS} 宏发股份2Q26收入59.14亿元，同比+36\%，毛利率33\%，符合预期。全年预计增长27\%（过去十年CAGR 15\%），其中约10个百分点来自提价。毛利率预计3Q26回升至34.5\%，4Q26至35.1\%。六大产品线均两位数增长，HVDC继电器+30\%，通用继电器+36\%。欧姆龙剥离继电器业务为宏发提供进一步扩大份额的窗口。
{\small\color{refgray}数据：2Q26收入59.14亿元，同比+36\%\par}
{\small\color{refgray}数据：全年预计增长27\%，提价贡献10个百分点\par}
{\small\color{refgray}数据：毛利率预计3Q26回升至34.5\%\par}
{\small\color{refgray}数据：欧姆龙剥离继电器业务\par}
{\small\color{refgray}边际变化：增长逻辑从跟随行业增长转变为主动夺取份额，定价能力显著增强。\par}
\paragraph{GS} 日立第一财季能源业务订单同比+87\%（剔除核电+94\%），基础订单同样强劲。订单积压10.3万亿日元，环比+12\%。公司上调全年指引，能源板块贡献绝大部分上调。管理层认为利润率仍有持续改善空间：订单积压利润率结构优于过去、AI应用提升效率、项目管理改进。AI利用率目标FY3/28达到30\%（当前约10\%）。
{\small\color{refgray}数据：能源订单同比+87\%，剔除核电+94\%\par}
{\small\color{refgray}数据：订单积压10.3万亿日元，环比+12\%\par}
{\small\color{refgray}数据：AI利用率目标FY3/28达30\%\par}
{\small\color{refgray}数据：上调全年销售与利润指引\par}
{\small\color{refgray}边际变化：能源业务订单增长并非仅靠大单拉动，基础订单广泛走强意味着需求具有持续性。\par}
\paragraph{GS} 保时捷2Q26集团EBIT 7.53亿欧元，超市场共识7.13亿，利润率8.5\%。超预期来自911高配车型（GTS/Turbo/GT）组合驱动，单车收入12.6万欧元。汽车业务净现金流5.06亿欧元，现金流利润率6.5\%超过全年指引3-5\%上限。研发支出低于预期，更多成本计入当期损益，传递管理层对盈利能力的信心。
{\small\color{refgray}数据：集团EBIT 7.53亿欧元，超共识7.13亿\par}
{\small\color{refgray}数据：单车收入12.6万欧元，高于预期12.5万\par}
{\small\color{refgray}数据：现金流利润率6.5\%，超全年指引上限\par}
{\small\color{refgray}数据：研发支出5.17亿欧元，低于预期5.96亿\par}
{\small\color{refgray}边际变化：盈利质量结构性改善，高配车型组合驱动，而非一次性项目。\par}
\paragraph{GS} 力拓上半年基础EBITDA 148亿美元，同比+28\%，净利润69亿美元，同比+43\%，超GS预期11\%。有效税率25\%，低于此前指引30\%。22个主要开发项目全部按计划或超前推进，包括西芒杜、奥尤陶勒盖等。成本削减目标从6.5亿提升至18亿美元（约占成本基础5\%），上半年已实现13亿年化。非核心资产出售计划50-100亿美元。
{\small\color{refgray}数据：基础EBITDA 148亿美元，同比+28\%\par}
{\small\color{refgray}数据：净利润69亿美元，同比+43\%\par}
{\small\color{refgray}数据：成本削减目标18亿美元，上半年已实现13亿年化\par}
{\small\color{refgray}数据：22个主要开发项目全部按计划或超前推进\par}
{\small\color{refgray}边际变化：成本削减力度比市场预期更激进，目标翻近3倍，显示管理层对运营效率的重视程度提升。\par}
\paragraph{GS} 施耐德电气2Q26有机销售增长16.5\%，超预期近6个百分点。数据中心业务三位数增长，但住宅业务在欧洲也出现环比改善。中国工业自动化增速19.6\%，远超预期7.7\%，与西门子形成正向共振。全年指引上调：有机销售增速从7-10\%提升至10-13\%，利润率目标从19.1-19.4\%上调至19.4-19.7\%。自由现金流超预期14\%。
{\small\color{refgray}数据：有机销售增长16.5\%，超预期10.6\%\par}
{\small\color{refgray}数据：中国工业自动化增速19.6\%，远超预期7.7\%\par}
{\small\color{refgray}数据：全年指引上调至10-13\%增长\par}
{\small\color{refgray}数据：自由现金流超预期14\%\par}
{\small\color{refgray}边际变化：增长不仅依赖数据中心，住宅业务环比改善，中国工业自动化可能预示制造业资本开支边际回暖。\par}
\paragraph{MS} Bloom Energy首次实现单季营收超10亿美元（10.65亿），超市场预期29\%。连续第二个季度上调全年指引：营收中位数上调约4.5亿美元，利润指引上调26\%。调整后毛利率34.3\%，超预期350bps，运营利润2.4亿美元，超预期73\%。经营现金流2.26亿美元，市场预期为负值。管理层对钪供应、项目延迟和产能瓶颈的担忧给出增量安抚。
{\small\color{refgray}数据：单季营收10.65亿美元，超预期29\%\par}
{\small\color{refgray}数据：全年营收指引上调约4.5亿美元\par}
{\small\color{refgray}数据：调整后毛利率34.3\%，超预期350bps\par}
{\small\color{refgray}数据：经营现金流2.26亿美元，市场预期-0.42亿\par}
{\small\color{refgray}边际变化：连续两个季度主动上调指引，管理层对订单可见度和交付节奏有较高信心。\par}
\subsubsection*{关键数据}
\begin{itemize}
  \item 中国6月变压器对美出口同比+172\%，环比+52\%
  \item 美国变压器供需缺口72\%，预计2028年收窄至57\%
  \item 全球变压器交货周期128周，中国供应商24-36周
  \item 巴斯夫现金转换率从89\%骤降至21\%
  \item 施耐德中国工业自动化增速19.6\%，远超预期7.7\%
  \item 高端CCL台光毛利率环比+4.4ppt至33.9\%
  \item Bloom Energy首次单季营收超10亿美元
  \item 力拓成本削减目标从6.5亿提升至18亿美元
  \item 宏发股份全年预计增长27\%，提价贡献10个百分点
  \item 日立能源订单同比+87\%，订单积压10.3万亿日元
\end{itemize}
\begin{figure}[htbp]
\centering
\includegraphics[width=0.88\linewidth]{figures/fig_028.jpg}
\caption{EXHIBIT 3 - EXHIBIT 3: HON orders acceleration \textbackslash{}- Double-digit short-cycle orders growth in all segments, driving total orders up \$16\textbackslash{}\%\$ organically \textbackslash{}- PA\&T orders up 24\% organically driven by surging LNG and gas processing demand globally}
\end{figure}
\begin{figure}[htbp]
\centering
\includegraphics[width=0.88\linewidth]{figures/fig_082.jpg}
\caption{Exhibit 1 - Exhibit 2: China total transformer export value picked up to \$47\textbackslash{}\%\$ yoy in June (vs \$72\textbackslash{}\%\$ in May)}
\end{figure}
\begin{figure}[htbp]
\centering
\includegraphics[width=0.88\linewidth]{figures/fig_090.jpg}
\caption{Exhibit 2 - Exhibit 3: We expect both EMC and TUC GM to start the next round of expansion driven by solid AI demand and tightening S/D outlook Exhibit 4: Both EMC and TUC will expand more than \$100\textbackslash{}\%+\$ capacity in coming 2 years, while both}
\end{figure}
\begin{figure}[htbp]
\centering
\includegraphics[width=0.88\linewidth]{figures/fig_104.jpg}
\caption{Exhibit 1 - Exhibit 1: Schneider 2Q26 results vs. company-compiled consensus Exhibit 2: Schneider Group OSG vs peers Exhibit 3: Schneider Energy Management OSG vs peers We are Buy rated on Schneider (also on the Conviction List), with a 12}
\end{figure}
{\small\color{refgray}本节综合 11 篇研究；涉及 Bernstein、NOM、GS、MS。}
\newpage
\subsection{AI变现进入验证期：广告效率与Copilot先行，消费者AI与主题乐园待考}
\textbf{2026年Q2财报季显示，AI在互联网与传媒领域的变现路径出现分化：Meta和微软的广告与Copilot业务已产生可量化回报，但消费者AI的粘性差异和主题乐园客流增速放缓构成主要不确定性。}
\subsubsection*{主流共识}
\begin{itemize}
  \item AI在核心广告业务中的效率提升已可量化，Meta和微软均报告了广告点击率、转化率或收入增速的改善。
  \item 企业级AI（如M365 Copilot、Azure）的付费采用正在加速，微软Copilot席位单季净增超预期。
  \item 大型科技公司资本支出维持高位，但短期收入增长为AI投入提供了合理性。
\end{itemize}
\subsubsection*{分歧与条件差异}
\begin{itemize}
  \item 消费者AI的变现路径存在分歧：Meta押注个人超级智能，而通义千问等产品用户粘性弱，流量质量存疑。
  \item 主题乐园客流增速：Citi认为美国本土客流增长放缓，而市场对Q4的6\%增长预期存在不确定性。
  \item AI代理的变现模式：Meta强调订阅、按使用付费等可能性，但具体路径尚不明确；微软则已通过Copilot实现席位增长。
\end{itemize}
\subsubsection*{机构观点}
\paragraph{NOM} 中国移动互联网流量向企业AI代理平台转移。豆包DAU达1.67亿，环比增6\%，同比增3.3倍，领先优势扩大。通义千问DAU达2930万，但55\%用户日均使用不足1分钟，粘性弱于豆包（85\%重度用户）和DeepSeek（81\%）。腾讯元宝月访问量从3月885万升至2100万，字节TRAE IDE以1279万次居次。
{\small\color{refgray}数据：豆包DAU 1.67亿，环比+6\%，同比+3.3倍\par}
{\small\color{refgray}数据：通义千问DAU 2930万，55\%用户日均使用<1分钟\par}
{\small\color{refgray}数据：腾讯元宝月访问量2100万，较3月+137\%\par}
{\small\color{refgray}边际变化：流量从消费级AI聊天机器人向企业级AI代理平台转移，通义千问用户增长快但粘性恶化。\par}
\paragraph{Bernstein} Meta 2Q26营收增长28\%，广告收入增长27\%，Advantage+年化收入超750亿美元。AI驱动Facebook广告点击量+8.3\%，转化率+15.7\%。但核心关注点是消费者AI战略：扎克伯格押注个人超级智能，以应对TikTok式冲击，36亿DAU是独特筹码。
{\small\color{refgray}数据：广告收入同比+27\%至594亿美元\par}
{\small\color{refgray}数据：Advantage+年化收入>750亿美元\par}
{\small\color{refgray}数据：Facebook广告点击量+8.3\%，转化率+15.7\%\par}
{\small\color{refgray}边际变化：AI投入在广告业务已产生可量化回报，但消费者AI战略昂贵且不确定，与谷歌、OpenAI的企业市场路径分化。\par}
\paragraph{Citi} 迪士尼FY26调整后EPS 16\%增长指引存在不确定性。预计FY26 EPS为6.80美元，略低于共识。第三财季EBIT预计52.35亿美元，低于公司指引53亿美元。主题乐园EBIT 29.54亿美元，同比+17.4\%，但美国本土客流增速仅1\%，Q4市场预期的6\%增长存疑。流媒体EBIT 6.1亿美元，同比+85.5\%，但增速低于预期。
{\small\color{refgray}数据：FY26调整后EPS预测6.80美元，同比+15\%\par}
{\small\color{refgray}数据：第三财季EBIT预测52.35亿美元，低于指引53亿美元\par}
{\small\color{refgray}数据：主题乐园EBIT 29.54亿美元，同比+17.4\%\par}
{\small\color{refgray}边际变化：市场焦点从单一增长叙事转向多业务线调整，主题乐园客流增速放缓，流媒体内容支出上升。\par}
\paragraph{GS} Meta 2Q26广告收入同比+28\%，Advantage+年化收入750亿美元。AI驱动的推荐和广告效率提升已可量化，但管理层对2027-2028年资本支出路径缺乏具体框架，GS上调累计资本支出预测约14\%。AI代理被视为核心业务之外的延伸，但变现路径尚不明确。
{\small\color{refgray}数据：广告收入同比+28\%\par}
{\small\color{refgray}数据：Advantage+年化收入750亿美元\par}
{\small\color{refgray}数据：2027-2028年累计资本支出预测上调约14\%\par}
{\small\color{refgray}边际变化：AI投入短期合理性得到收入增长支撑，但长期资本支出框架仍为定性描述，市场疑虑未消。\par}
\paragraph{MS} 微软F4Q26 Azure增速43\%，超指引上限3个百分点，F1Q27指引约45\%。M365 Copilot付费席位突破3000万，单季净增约1000万，超市场预期。运营利润率45.1\%，超预期70个基点。ARPU增长拆解为三个引擎：Copilot席位、E7迁移、消费型变现。
{\small\color{refgray}数据：Azure固定汇率增长43\%，超指引上限3个百分点\par}
{\small\color{refgray}数据：M365 Copilot付费席位>3000万，单季净增约1000万\par}
{\small\color{refgray}数据：运营利润率45.1\%，超预期70个基点\par}
{\small\color{refgray}边际变化：AI变现从“能否发生”推进到“能持续多久”，Azure增长瓶颈在供给端而非需求端。\par}
\paragraph{MS} Meta核心社交业务用户规模与变现效率同步提升。应用家族DAU达36亿，Instagram单平台DAU突破20亿。第三季度收入指引上限640亿美元，隐含同比+26\%。四项新产品期权（Neocloud、API、订阅、消息商业化）合计可为2028年EPS贡献约9美元（约25\%上行空间）。
{\small\color{refgray}数据：应用家族DAU 36亿，Instagram DAU 20亿\par}
{\small\color{refgray}数据：第三季度收入指引上限640亿美元，同比+26\%\par}
{\small\color{refgray}数据：新产品期权合计可为2028年EPS贡献约9美元\par}
{\small\color{refgray}边际变化：AI+大模型使广告匹配更精准，多代际护城河加深；新产品期权尚未被市场充分定价。\par}
\subsubsection*{关键数据}
\begin{itemize}
  \item 豆包DAU 1.67亿，同比+3.3倍
  \item 通义千问55\%用户日均使用<1分钟
  \item Meta广告收入同比+27\%至594亿美元
  \item Advantage+年化收入>750亿美元
  \item 微软Azure增速43\%，超指引上限3个百分点
  \item M365 Copilot付费席位>3000万，单季净增约1000万
  \item 迪士尼FY26 EPS预测6.80美元，同比+15\%
  \item Meta第三季度收入指引上限640亿美元，同比+26\%
  \item Meta新产品期权合计可为2028年EPS贡献约9美元
\end{itemize}
\begin{figure}[htbp]
\centering
\includegraphics[width=0.88\linewidth]{figures/fig_059.jpg}
\caption{EXHIBIT 2 - EXHIBIT 2: Meta 2Q26: EBIT is impacted by one-time legal charge EXHIBIT 3: Digital ads growth (\%) Reported ad revenue Y/Y growth (\%) through 2Q26 EXHIBIT 4: Revenue contribution by geography}
\end{figure}
\begin{figure}[htbp]
\centering
\includegraphics[width=0.88\linewidth]{figures/fig_074.jpg}
\caption{Figure 5 - \#\# \#2: US Attendance Trends Over the past few quarters, investors have closely followed Domestic attendance. To gauge recent trends, we looked at three external data sources: ■ First, in F3Q26, Orlando airport passengers (through May) grew 3\% YoY. Figure 5. Or}
\end{figure}
\begin{figure}[htbp]
\centering
\includegraphics[width=0.88\linewidth]{figures/fig_137.jpg}
\caption{Exhibit 1 - Exhibit 1: Model Changes}
\end{figure}
\begin{figure}[htbp]
\centering
\includegraphics[width=0.88\linewidth]{figures/fig_138.jpg}
\caption{Exhibit 1 - Exhibit 2: Our \textbackslash{}\$775 PT Implies Paying 23X our average '27/'28 EPS of \textbackslash{}\$34/\textbackslash{}\$35, representing 40\% upside}
\end{figure}
{\small\color{refgray}本节综合 6 篇研究；涉及 NOM、Bernstein、Citi、GS、MS。}
\newpage
\subsection{医疗与保险：分化加剧，创新与资本效率成关键}
\textbf{中国医疗与保险市场正经历结构性分化：跨国药企在华业绩因产品组合与政策暴露度不同而显著分化，保险行业则因资本缓冲差距加速两极分化。创新产品准入和成本控制成为区分竞争力的核心变量。}
\subsubsection*{主流共识}
\begin{itemize}
  \item 跨国药企在华业务受医保谈判和集采政策影响，业绩表现分化，创新产品（如诺华Leqvio）成为增长亮点。
  \item 中国保险行业集中度持续提升，头部寿险和财险公司凭借资本优势和品牌认知占据主导地位。
  \item GSK等药企通过成本节约计划为研发和业务拓展提供资金，但市场已部分预期其利润率改善。
\end{itemize}
\subsubsection*{分歧与条件差异}
\begin{itemize}
  \item 跨国药企在华业绩：罗氏药品销售同比下滑33\%，而诺华增长12\%，差异源于产品组合对医保政策的敏感度不同。
  \item 保险资本缓冲：上市寿险公司核心偿付能力普遍高于100\%，而17家未上市险企低于100\%，其中4家低于70\%，逼近监管红线。
  \item GSK成本节约的增量贡献：市场共识已包含2026-2030年约40个基点的利润率提升，GSK额外成本节约仅能再贡献约20个基点，对2030年核心EBIT上调空间约1\%。
\end{itemize}
\subsubsection*{机构观点}
\paragraph{NOM} 跨国药企在华业绩分化显著：罗氏中国区药品销售同比下滑33\%至5.22亿瑞士法郎，诊断业务下滑15\%，受医疗定价改革影响；阿斯利康中国区销售下降7\%（剔除汇率影响下降13\%）至15.9亿美元，集采持续冲击；诺华中国区增长12\%（剔除汇率影响增长6\%）至12亿美元，Leqvio纳入医保后使用量快速提升。创新产品准入和医保谈判结果成为区分业绩的关键。
{\small\color{refgray}数据：罗氏中国区药品销售同比下滑33\%至5.22亿瑞士法郎\par}
{\small\color{refgray}数据：阿斯利康中国区销售下降7\%（剔除汇率影响下降13\%）至15.9亿美元\par}
{\small\color{refgray}数据：诺华中国区增长12\%（剔除汇率影响增长6\%）至12亿美元\par}
{\small\color{refgray}数据：罗氏诊断业务同比下滑15\%至3.9亿瑞士法郎\par}
{\small\color{refgray}边际变化：跨国药企在华业务从统一增长转向分化，诺华因Leqvio纳入医保实现逆势增长，而罗氏和阿斯利康受政策冲击更大。\par}
\paragraph{JPM} 中国保险行业两极分化加速：头部五家寿险和财险公司分别占据45\%和68\%市场份额，而中小险企资本缓冲不足。截至2026年3月，17家未上市寿险公司核心偿付能力低于100\%，其中4家低于70\%，逼近50\%监管红线。2025年55\%的财险公司出现承保亏损。规模成为区分竞争力的关键变量，中小险企在品牌、渠道和资本方面均处劣势。
{\small\color{refgray}数据：前5大寿险公司市场份额45\%，前5大财险公司市场份额68\%\par}
{\small\color{refgray}数据：17家未上市寿险公司核心偿付能力低于100\%，其中4家低于70\%\par}
{\small\color{refgray}数据：2025年55\%的财险公司出现承保亏损\par}
{\small\color{refgray}数据：保险深度从2015年3.5\%升至2025年4.4\%\par}
{\small\color{refgray}边际变化：保险行业从数量扩张转向质量分化，资本缓冲差距成为区分上市险企与中小险企的核心指标，中小险企面临资本补充压力。\par}
\paragraph{JPM} GSK重申2031年销售额超400亿英镑目标，并宣布到2029年实现19亿英镑年化成本节约的重组计划。JPM测算成本节约对核心EBIT净贡献约2亿英镑（约60个基点利润率提升），但市场共识已包含2026-2030年约40个基点的利润率改善，额外成本节约仅能再贡献约20个基点，对应2030年核心EBIT约1\%上调空间。关键三期临床数据集中在2028-2030年读出，管线进展才是长期预测上调的真正驱动。
{\small\color{refgray}数据：GSK 2031年销售目标超400亿英镑\par}
{\small\color{refgray}数据：重组计划到2029年实现19亿英镑年化成本节约\par}
{\small\color{refgray}数据：成本节约对核心EBIT净贡献约2亿英镑（约60个基点）\par}
{\small\color{refgray}数据：市场共识已包含2026-2030年约40个基点的利润率提升\par}
{\small\color{refgray}边际变化：GSK成本节约计划已被市场部分预期，增量贡献有限，市场关注点转向2028-2030年关键管线数据（如ris-rez、efimosfermin、velzatinib）。\par}
\subsubsection*{关键数据}
\begin{itemize}
  \item 罗氏中国区药品销售同比下滑33\%至5.22亿瑞士法郎
  \item 诺华中国区增长12\%（剔除汇率影响增长6\%）至12亿美元
  \item 17家未上市寿险公司核心偿付能力低于100\%，其中4家低于70\%
  \item 2025年55\%的财险公司出现承保亏损
  \item GSK成本节约对核心EBIT净贡献约2亿英镑，增量利润率约20个基点
\end{itemize}
\begin{figure}[htbp]
\centering
\includegraphics[width=0.88\linewidth]{figures/fig_106.jpg}
\caption{Figure 1 - Figure 1: China – time series data for licensed insurers Number of legal entities Number of legal entities}
\end{figure}
\begin{figure}[htbp]
\centering
\includegraphics[width=0.88\linewidth]{figures/fig_107.jpg}
\caption{Figure 2 - Figure 2: China Insurance – life insurers and P\&C insurers' premium trend Rmb in billions, \%}
\end{figure}
\begin{figure}[htbp]
\centering
\includegraphics[width=0.88\linewidth]{figures/fig_108.jpg}
\caption{Figure 3 - Figure 3: China life insurers – core solvency ratio distribution \% Figure 4: China non-life insurers – core solvency ratio distribution \#\# Risk management score and insurance risk rating}
\end{figure}
\begin{figure}[htbp]
\centering
\includegraphics[width=0.88\linewidth]{figures/fig_109.jpg}
\caption{Figure 3 - Figure 3: China life insurers – core solvency ratio distribution \% Figure 4: China non-life insurers – core solvency ratio distribution \#\# Risk management score and insurance risk rating We have illustrated the divergence in c}
\end{figure}
{\small\color{refgray}本节综合 3 篇研究；涉及 NOM、JPM。}
\newpage
\subsection{区域市场：亚洲分化与科技去杠杆尾声}
\textbf{亚洲市场正经历结构性分化：中国智慧城市与其余地区形成K型增长，科技板块去杠杆接近尾声，而日本汽车受益于美国稳定定价，中国汽车等待新周期。}
\subsubsection*{主流共识}
\begin{itemize}
  \item 中国智慧城市（7城）贡献了全国上半年20\%的增长，为二十年最高，AI驱动增长集中化。
  \item 全球科技与半导体板块去杠杆速度超预期，进一步空间有限，市场接近阶段性底部。
  \item 美国汽车定价环境稳定，福特上调全年EBIT指引，利好高美国敞口的日本车企。
  \item 中国汽车新车发布效应消退，多数车企订单环比回落，市场等待下一产品周期。
\end{itemize}
\subsubsection*{分歧与条件差异}
\begin{itemize}
  \item 中国智慧城市增速加速（5.4\%→5.6\%），而其余83\%地区增速放缓（4.9\%→4.5\%），K型分化加剧。
  \item 人民币中间价模型显示逆周期因子修正约191个基点，政策意图与市场定价存在张力。
  \item 日本汽车中，本田、丰田、斯巴鲁因混动组合和美国敞口被看好，但贸易政策（USMCA）带来不确定性。
  \item 中国车企中，理想和小鹏因新车型订单环比正增长，而比亚迪、吉利等处于产品空窗期，订单偏弱。
\end{itemize}
\subsubsection*{机构观点}
\paragraph{NOM} 中国宏观增长出现地理K型分化：2026年上半年全国GDP增速从5.0\%降至4.7\%，但七个智慧城市（3个一线+4个精选二线）加权平均增速从5.4\%升至5.6\%，贡献全国增长的20\%，为至少二十年最高。其余占GDP 83\%的地区增速从4.9\%降至4.5\%。AI产业集聚效应是核心驱动力，北京和上海固定资产投资分别增长3.0\%和6.8\%，而全国下降5.7\%。这一分化非短期波动，全面刺激可能性降低。
{\small\color{refgray}数据：全国GDP增速从5.0\%降至4.7\%\par}
{\small\color{refgray}数据：智慧城市加权平均增速从5.4\%升至5.6\%\par}
{\small\color{refgray}数据：智慧城市贡献全国增长的20\%，二十年最高\par}
{\small\color{refgray}数据：其余地区增速从4.9\%降至4.5\%\par}
{\small\color{refgray}边际变化：此前市场关注总量增速放缓，但NOM指出AI红利仅被少数智慧城市捕获，其余地区增速下滑幅度（0.4个百分点）是总量降幅（0.3个百分点）的1.3倍，K型分化在加速。\par}
\paragraph{NOM} 美元兑人民币中间价模型预测值下调至6.7536，较前次低约119个基点；引入逆周期因子后预测值为6.7727，较前次中间价低约172个基点。逆周期因子对模型修正约191个基点，显示政策层面仍在施加缓冲。美元指数变动是最大负向贡献因子，模型误差近期呈现系统性偏大特征，表明输入变量与实际定价存在结构性错位。
{\small\color{refgray}数据：模型预测值6.7536，较前次低119个基点\par}
{\small\color{refgray}数据：引入逆周期因子后预测值6.7727，较前次低172个基点\par}
{\small\color{refgray}数据：逆周期因子修正幅度约191个基点\par}
{\small\color{refgray}数据：美元指数变动是最大负向贡献因子\par}
{\small\color{refgray}边际变化：此前模型误差为随机扰动，近期连续同向偏离且幅度放大，表明模型依赖的利率、汇率波动率等变量与实际市场定价出现结构性错位，而非单纯随机波动。\par}
\paragraph{GS} 福特Q2财报上调全年调整后EBIT指引至100-110亿美元（此前85-105亿美元），主要来自Ford Blue（燃油/混动业务）EBIT展望上调至50-55亿美元。美国市场定价和产品组合稳定，全年定价预计同比增长0.5\%，皮卡需求强劲，V8发动机占比上升。储能系统订单积累，目标2028年实现20GWh产能。维持对本田、丰田、斯巴鲁的积极看法，因其高美国敞口和强混动组合。
{\small\color{refgray}数据：福特全年调整后EBIT指引上调至100-110亿美元\par}
{\small\color{refgray}数据：Ford Blue全年EBIT展望上调至50-55亿美元\par}
{\small\color{refgray}数据：美国市场全年定价预计同比增长0.5\%\par}
{\small\color{refgray}数据：储能系统目标2028年20GWh产能\par}
{\small\color{refgray}边际变化：此前市场担忧美国汽车定价环境恶化，但福特财报显示定价稳定且混动需求强劲，日本车企凭借混动差异化优势可能获得更有利定价空间。\par}
\paragraph{JPM} 全球科技与半导体板块去杠杆速度超预期，进一步空间有限。对冲基金在SOX指数上的杠杆头寸已基本解除；CTA等动量交易者在纳指、KOSPI、台湾加权、日经上的多仓大幅削减；杠杆ETF的AUM占市值比已恢复85\%（半导体非存储）和55\%（存储），不再是主要波动源；散户看涨期权买量从6月5日峰值1400万张回落至历史均值附近；风险平价杠杆指标正常化。唯一未确认的是保证金账户杠杆（7月数据8月底公布）。
{\small\color{refgray}数据：杠杆ETF AUM占市值比恢复85\%（半导体非存储）和55\%（存储）\par}
{\small\color{refgray}数据：散户看涨期权买量从1400万张峰值回落至历史均值\par}
{\small\color{refgray}数据：对冲基金SOX杠杆头寸基本解除\par}
{\small\color{refgray}数据：CTA多仓在纳指、KOSPI、台湾加权、日经大幅削减\par}
{\small\color{refgray}边际变化：此前JPM预期去杠杆将持续更久，但实际推进速度更快，目前进一步空间有限，市场接近阶段性底部。杠杆ETF的凸性自修正机制加速了去杠杆进程。\par}
\paragraph{MS} 7月第三周中国主要车企订单分化：9家车企周度订单环比下降，比亚迪周订单7.05-7.1万辆（环比-8\%，同比-7\%），吉利银河、蔚来、零跑等均下滑。例外是理想（周订单9.2-9.4万辆，环比-49\%但7月累计环比+3\%，同比+18\%）和小鹏（周订单9.9-10.1万辆，环比-76\%但7月累计环比+43\%，同比+13\%），新车型L6和MONA M03仍在提供增量。市场进入等待下一产品周期（比亚迪“大汉”、吉利新车）的过渡阶段。
{\small\color{refgray}数据：比亚迪周订单7.05-7.1万辆，环比-8\%，同比-7\%\par}
{\small\color{refgray}数据：理想周订单9.2-9.4万辆，7月累计环比+3\%，同比+18\%\par}
{\small\color{refgray}数据：小鹏周订单9.9-10.1万辆，7月累计环比+43\%，同比+13\%\par}
{\small\color{refgray}数据：蔚来周订单1.15-1.2万辆，同比+121\%\par}
{\small\color{refgray}边际变化：此前7月初新车发布带来订单脉冲，但第三周效应消退，多数车企环比回落。理想和小鹏的月度环比仍为正，表明新车型效应周期缩短，但相对水平高于6月。\par}
\subsubsection*{关键数据}
\begin{itemize}
  \item 中国智慧城市贡献全国上半年增长的20\%，二十年最高
  \item 全国GDP增速从5.0\%降至4.7\%，智慧城市从5.4\%升至5.6\%
  \item 逆周期因子对人民币中间价模型修正约191个基点
  \item 福特全年EBIT指引上调至100-110亿美元
  \item 杠杆ETF AUM占市值比恢复85\%（半导体非存储）和55\%（存储）
  \item 散户看涨期权买量从1400万张峰值回落至历史均值
  \item 比亚迪周订单7.05-7.1万辆，环比-8\%
  \item 理想7月累计订单环比+3\%，小鹏环比+43\%
\end{itemize}
\begin{figure}[htbp]
\centering
\includegraphics[width=0.88\linewidth]{figures/fig_121.jpg}
\caption{Figure 1 - Figure 2: Short interest on Semiconductors ex Memory \& Hyperscalers individual stocks Short Interest as a \% share of shares outstanding.}
\end{figure}
\begin{figure}[htbp]
\centering
\includegraphics[width=0.88\linewidth]{figures/fig_123.jpg}
\caption{Figure 2 - Figure 4: Momentum signals for major equity indices}
\end{figure}
\begin{figure}[htbp]
\centering
\includegraphics[width=0.88\linewidth]{figures/fig_124.jpg}
\caption{Figure 4 - Figure 4: Momentum signals for major equity indices Figure 5: Momentum signals for Asian equity indices weeks. As a result, the reduction in the AUM of leveraged equity ETFs has taken place much faster than we previously envisag}
\end{figure}
\begin{figure}[htbp]
\centering
\includegraphics[width=0.88\linewidth]{figures/fig_122.jpg}
\caption{Figure 1 - Figure 2: Short interest on Semiconductors ex Memory \& Hyperscalers individual stocks Short Interest as a \% share of shares outstanding. Figure 3: Short interest on SMH and DRAM US Equity ETFs Short Interest as a \% share of shar}
\end{figure}
{\small\color{refgray}本节综合 5 篇研究；涉及 NOM、GS、JPM、MS。}
\newpage
\subsection{科技硬件与半导体设备：需求强劲，产能约束与解禁压力并存}
\textbf{半导体设备需求在AI和存储升级驱动下持续走强，Lam Research业绩与指引全面超预期，毛利率创20年新高；Seagate近线HDD产能已锁定至2028年，定价能力增强。但SpaceX面临8月大规模解禁压力，市场情绪谨慎，短期报价承压。}
\subsubsection*{主流共识}
\begin{itemize}
  \item AI和云服务商对高容量存储的长期需求正在重塑采购模式，Seagate近线HDD产能已分配至2028年。
  \item 半导体设备需求强劲，Lam Research下季度营收指引81亿美元，远超市场预期。
  \item 存储和半导体设备行业定价能力增强，毛利率和价格同比上涨。
\end{itemize}
\subsubsection*{分歧与条件差异}
\begin{itemize}
  \item Lam Research中国收入占比下降（从34\%降至26\%），但被NAND和逻辑业务增长抵消；市场对后续中国需求影响存在分歧。
  \item SpaceX短期报价受技术性因素（解禁）主导，而非基本面；部分投资者认为估值已反映AI业务价值为负，但多数持谨慎态度。
\end{itemize}
\subsubsection*{机构观点}
\paragraph{Bernstein} Lam Research 6月季度营收67.2亿美元、每股收益1.82美元，均超预期，下季度营收指引81亿美元（市场共识71.3亿）。毛利率52\%创20年新高，受益于定价、运营效率和产品组合。NAND业务环比增长118\%，逻辑增长63\%，中国收入占比从34\%降至26\%。2026年WFE预期上调至低1500亿美元区间，先进封装收入预计增长超70\%。
{\small\color{refgray}数据：营收67.2亿美元（预期66.8亿）\par}
{\small\color{refgray}数据：每股收益1.82美元（预期1.70）\par}
{\small\color{refgray}数据：毛利率52\%（预期50.5\%）\par}
{\small\color{refgray}数据：NAND业务环比增长118\%\par}
{\small\color{refgray}边际变化：毛利率突破52\%并上调长期目标至55\%，反映定价能力和产品结构发生结构性变化；WFE预期上调显示行业需求持续走强。\par}
\paragraph{MS} SpaceX二季度收入预计约67.5亿美元（略低于共识69.2亿），调整后EBITDA约20.2亿美元。真正压力来自8月6日首批9.3亿股解禁（约1000亿美元市值），总计近40亿股将在2027年1月前解锁。报价自高点下跌近50\%，低于IPO发行价。超过2/3的读者对年底前表现持谨慎态度，首个正面催化剂可能是8月底Starship第14次试飞。
{\small\color{refgray}数据：二季度收入预计67.5亿美元（共识69.2亿）\par}
{\small\color{refgray}数据：调整后EBITDA预计20.2亿美元（共识21.1亿）\par}
{\small\color{refgray}数据：8月6日首批9.3亿股解禁（约1000亿美元）\par}
{\small\color{refgray}数据：总计近40亿股将在2027年1月前解锁\par}
{\small\color{refgray}边际变化：市场关注点从季度业绩转向解禁压力，报价已低于IPO发行价，估值反映AI业务价值为负或零。\par}
\paragraph{NOM} Seagate 26/6财年Q4营收36.29亿美元（超预期34.9亿），营业利润率44.6\%（超预期41.9\%），Q1指引营收中值41亿美元、营业利润率约50\%。近线HDD 2028年供应已全部被客户锁定，客户已开始讨论2029年合同。每Exabyte价格同比上涨超10\%，公司通过技术迁移而非增加驱动器数量来扩产。
{\small\color{refgray}数据：Q4营收36.29亿美元（预期34.9亿）\par}
{\small\color{refgray}数据：营业利润率44.6\%（预期41.9\%）\par}
{\small\color{refgray}数据：Q1指引营收中值41亿美元\par}
{\small\color{refgray}数据：Q1指引营业利润率约50\%\par}
{\small\color{refgray}边际变化：客户采购周期从12-18个月延长至2028年，并开始2029年合同谈判，反映AI和云服务商对高容量存储的长期需求重塑采购模式。\par}
\subsubsection*{关键数据}
\begin{itemize}
  \item Lam Research营收67.2亿美元，超预期66.8亿
  \item 毛利率52\%创20年新高
  \item NAND业务环比增长118\%
  \item 中国收入占比从34\%降至26\%
  \item 2026年WFE预期上调至低1500亿美元
  \item SpaceX首批9.3亿股解禁（约1000亿美元）
  \item Seagate Q4营收36.29亿美元，超预期34.9亿
  \item 近线HDD 2028年供应已全部被锁定
  \item 每Exabyte价格同比上涨超10\%
\end{itemize}
\begin{figure}[htbp]
\centering
\includegraphics[width=0.88\linewidth]{figures/fig_023.jpg}
\caption{EXHIBIT 2 - EXHIBIT 4: NAND and Logic saw strong growth QoQ, while Foundry and DRAM were down slightly QoQ EXHIBIT 5: Memory share was up 7ppt QoQ to 46\% with equal share from DRAM and NAND Lam Research}
\end{figure}
\begin{figure}[htbp]
\centering
\includegraphics[width=0.88\linewidth]{figures/fig_044.jpg}
\caption{Exhibit 2 - Exhibit 2: 2Q26 MS \& Consensus: Segment Detail and Key KPIs Exhibit 3: 2028 EV/EBIT: SpaceX vs. Key Comps Note: Market data as of 7/27/2026 Exhibit 4: Short Interest vs. Key Comps Note: Market data as of 7/27/2026}
\end{figure}
\begin{figure}[htbp]
\centering
\includegraphics[width=0.88\linewidth]{figures/fig_024.jpg}
\caption{EXHIBIT 3 - EXHIBIT 5: Memory share was up 7ppt QoQ to 46\% with equal share from DRAM and NAND Lam Research Memory Contribution to Systems Revenues}
\end{figure}
\begin{figure}[htbp]
\centering
\includegraphics[width=0.88\linewidth]{figures/fig_045.jpg}
\caption{Exhibit 3 - Exhibit 3: 2028 EV/EBIT: SpaceX vs. Key Comps Note: Market data as of 7/27/2026 Exhibit 4: Short Interest vs. Key Comps Note: Market data as of 7/27/2026 \#\# Valuation Methodology and Risks}
\end{figure}
{\small\color{refgray}本节综合 3 篇研究；涉及 Bernstein、MS、NOM。}
\newpage
\section{辅助信号｜战略咨询}
本板块整合波士顿咨询与麦肯锡最新研究，聚焦三大主题：航空业成本重构与利润分化、城市竞争逻辑向居民体验转变、以及AI部署从实验转向组织变革。各报告均被实质纳入对应主题。
\subsection{航空业成本重构与利润分化}
\textbf{全球航空业单位成本增速已全面超过单位收入增速，利润分布极不均匀，竞争焦点从恢复运力转向重构成本结构。}
\begin{itemize}
  \item 截至2025年Q3，全球所有地区单位成本（CASK，不含燃油）增速均超过单位收入（PRASK）增速；机组人员成本同比上涨5\%-7\%，地面处理成本上涨4\%-7\%，人工成本是核心通胀驱动因素。
  \item 全服务航司（FSC）与低成本航司（LCC）利润分化加剧：北美FSC靠辅助收入维持，欧洲LCC利润率承压，亚太LCC表现优于FSC。
  \item AI驱动的预测性维护和机组调度优化成为永久性成本改善的关键工具，而非临时性费用削减。
  \item OEM交付缓解和新机型（如A321XLR）将重塑航线网络，但成本曲线重塑能力决定航司下一阶段相对位置。
\end{itemize}
\subsection{城市竞争转向居民体验与变革速度}
\textbf{城市竞争逻辑已从基础设施资产转向居民体验价值，变革速度成为可持续优势的新来源，中东和中国城市主导全球变革最快阵营。}
\begin{itemize}
  \item BCG《Cities of Choice 4.0》报告引入“居民倡导度”指标，通过满意度、推荐意愿等五维度测量居民对城市的长期承诺，倡导度高的城市居民留居比例更高。
  \item 变革速度排名中，深圳（92分）、利雅得（91分）、塔什干（88分）位列前三；前十名中中国城市占四席（深圳、广州、上海、北京），沙特城市占四席。
  \item 线上活跃居民在超过70个城市中的倡导度高于其他群体，数字服务能力成为城市留住人才的关键杠杆。
  \item 欧洲城市变革速度整体偏慢，但苏黎世等城市在宜居性维度仍具优势，反映不同城市发展路径的分化。
\end{itemize}
\subsection{AI部署从实验转向组织变革}
\textbf{企业AI部署正从实验性转向系统性变革，工作流程重新设计和CEO直接参与治理是影响利润的最强因素，但企业级利润影响仍不显著。}
\begin{itemize}
  \item 麦肯锡2025年3月全球调研显示，21\%的受访企业因部署生成式AI而从根本上重新设计了至少部分工作流程，工作流程重新设计对EBIT影响最大。
  \item CEO直接监督AI治理的企业利润回报更高：28\%的受访企业由CEO负责AI治理，17\%由董事会监督，大型企业中CEO监督对EBIT影响最显著。
  \item 超过80\%的受访者称生成式AI尚未对企业级EBIT产生实质性影响，业务单元层面的成本削减尚未转化为企业级利润变化。
  \item 大型企业在AI治理、风险管理和规模化路线图方面领先，正在积累难以追赶的竞争优势。
\end{itemize}
\begin{figure}[htbp]
\centering
\includegraphics[width=0.88\linewidth]{figures/fig_155.jpg}
\caption{Exhibit 1 - The survey findings also shed light on how organizations are structuring their AI deployment efforts. Some essential elements for deploying AI tend to be fully or partially centralized (Exhibit 1). For risk and compliance, as well as data governance, organizat}
\end{figure}
\begin{figure}[htbp]
\centering
\includegraphics[width=0.88\linewidth]{figures/fig_156.jpg}
\caption{Exhibit 2 - Organizations have employees overseeing the quality of gen AI outputs, though the extent of that oversight varies widely. Twenty-seven percent of respondents whose organizations use gen AI say that employees review all content created by gen AI before it is us}
\end{figure}
\subsection{咨询与企业战略}
\textbf{用于补充企业战略、管理层议题和产业落地视角。}
\begin{itemize}
  \item 消费者支出节奏变化、通胀侵蚀购买力、关税波动带来不确定性——波士顿咨询在一份2026年1月发布的报告中指出，消费品企业正面临一个“残酷的环境”。但报告的核心判断并非描述困境，而是指出：多数企业正在采用的供应链效率提升、价格微调和促销调整等传统手段，已经不足以应对当前局面。真正能拉...
  \item 全球主权财富产品和公共养老产品的管理资产规模在2025年达到43万亿美元，较2015年的21万亿美元翻了一番以上。波士顿咨询在最新发布的《全球主要读者报告》中指出，这一资本池已占全球专业管理资产总额约30\%，其配置决策足以对市场流动性和市场定价产生系统性影响。报告同时观察到，主权...
  \item 波士顿咨询一项跨行业调查显示，约三分之一的合资企业失败至少部分可归因于董事会未能维持清晰的决策权、规范的治理流程以及以合资企业利益而非母公司利益为优先的承诺。这一问题正随着更多企业将合资作为创造新价值的路径而愈发突出。该机构认为，治理失效已成为合资企业失败的第三大常见原因，其影响...
\end{itemize}
\begin{figure}[htbp]
\centering
\includegraphics[width=0.88\linewidth]{figures/fig_154.jpg}
\caption{Exhibit 2 - \# Lessons for Stakeholders JV leaders and stakeholders can use the board model lessons (see Exhibit 2) to help understand their roles and embrace an appropriate governance mindset. \#\# EXHIBIT 2 \#\# Lessons from Three Board Models Heads of Legal, Strategy, and C}
\end{figure}
\newpage
\section{辅助信号｜智库/国际机构}
本板块整合亚洲开发银行、IMF、经合组织等机构最新研究，聚焦亚太数字贸易结构性分化、全球排放核算方法升级、以及新兴经济体税收透明度与半导体生态建设三大主题，提供对宏观、产业及资产叙事有辅助价值的信号。
\subsection{亚太数字贸易：增长强劲但红利分配不均，监管碎片化成主要障碍}
\textbf{亚太地区贡献全球近三分之二电子商务销售额，移动支付覆盖约70\%交易，但数字鸿沟、合规成本上升和信任缺失导致各国受益程度差异显著，政策制定者需构建可互操作的监管框架。}
\begin{itemize}
  \item 全球电商销售额2016-2022年增长近60\%至27万亿美元，增速最快10国中7个在亚太；亚太占全球电商近2/3，移动支付覆盖约70\%交易。
  \item 数字服务出口自2005年以来增长近四倍，社交电商、直播购物等新业态在年轻消费者中快速普及。
  \item 红利分布极不均衡：中国台北、韩国、新加坡数字宏观环境对GDP贡献率分别为6.1\%、5.8\%和5.4\%，而菲律宾仅2.5\%，老挝和柬埔寨仅1.8\%。
  \item 结构性障碍包括监管碎片化、合规成本上升和信任缺失；约1/3人口仍未接入互联网，数字鸿沟叠加基础设施、数字素养和制度能力差距。
  \item 印尼电子商务已占其数字宏观环境72\%，预计2025年达1460亿美元，占GDP比重持续上升，反映新兴市场仍有显著增长空间。
\end{itemize}
\begin{figure}[htbp]
\centering
\includegraphics[width=0.88\linewidth]{figures/fig_139.jpg}
\caption{Figure 1 - Figure 1: Digitalization Index, 2024 This policy brief analyzes key trends, patterns, and impacts of digital trade and cross-border e-commerce in Asia and the Pacific and identifies regulatory challenges and opportunities. It dr}
\end{figure}
\begin{figure}[htbp]
\centering
\includegraphics[width=0.88\linewidth]{figures/fig_140.jpg}
\caption{Figure 2 - Figure 2: Trade in Digitally Delivered Services in Asia and the Pacific, 2010–2024 Trade agreements are critical in providing an enabling regulatory environment for cross-border e-commerce. Multilateral and regional frameworks h}
\end{figure}
\subsection{全球排放核算升级：居民原则调整与国际运输纳入将重塑国家间排放比较基准}
\textbf{经合组织新方法通过引入居民原则、扩展时间序列至1990年并首次纳入国际航空与海运排放，有望显著改善排放数据的政策可用性，但工业渔船排放可能被系统性低估。}
\begin{itemize}
  \item 全球仅53个国家在过去五年内编制过官方空气排放账户，各国在气体范围、时间序列长度和行业细分程度上差异巨大。
  \item 新方法从领土原则转向居民原则：扣除非居民在本国领土产生的排放，同时加入本国居民在境外产生的排放，使排放账户与国民宏观环境核算体系一致。
  \item 国际航空和海运排放首次被纳入核算边界，按居民归属重新分配；对于拥有大型国际航运或航空公司的国家，排放责任将显著上升。
  \item 行业细分扩展至64个（A*64标准），时间跨度从1990年至前年（t-2），数据来源从单一转向多源整合。
  \item 工业渔船排放可能被系统性低估，因公开自动识别系统数据无法覆盖大量远洋作业渔船，数据质量仍存在不确定性。
\end{itemize}
\begin{figure}[htbp]
\centering
\includegraphics[width=0.88\linewidth]{figures/fig_144.jpg}
\caption{Figure 1 - \#\# Preliminary results Figure 1. CO \$\_\{2\}\$ emissions with industry/household breakdown Unit: Million tonnes Note: ISIC A “Agriculture, forestry and fishing”, B “Mining and quarrying”, C “Manufacturing”, D “Electricity, gas, steam and air conditioning supply”, }
\end{figure}
\begin{figure}[htbp]
\centering
\includegraphics[width=0.88\linewidth]{figures/fig_146.jpg}
\caption{Figure 2 - 74. Figure 2 compares total emissions from the AEA estimates with national inventory totals (CRT and PRIMAP), as well as with energy-related emissions in the IEA database, for the two countries shown in Figure 1. Emission levels from CRT and PRIMAP are often v}
\end{figure}
\subsection{新兴经济体制度基建：税收透明度与半导体生态建设并行推进}
\textbf{坦桑尼亚、纳米比亚、库克群岛在税收信息交换上获'大体合规'评级，但受益所有权信息可用性存在结构性缺口；巴拿马试图利用区位优势发展半导体生态，但面临人才、基础设施和本地企业生态三重制约。}
\begin{itemize}
  \item 坦桑尼亚获全球论坛'大体合规'评级，但所有权与身份信息可用性仅获'部分合规'；法律允许名义持股和不记名公司权证，实际控制人可通过这两类安排隐藏身份。
  \item 纳米比亚2023年引入受益所有权登记制度，截至评审期末54\%公司、44\%封闭公司和63\%信托已提交信息，但改革时间太短，监督和执法仍处于萌芽状态。
  \item 库克群岛信息可用性三个核心要素评级均从'合规'下调至'大体合规'；国内合伙与信托不必然与反洗钱义务主体建立关系，名义持股无披露义务。
  \item 巴拿马高技术企业仅占制造业企业总数不到5\%，且高度集中在首都；2021-2022年企业固定资产总值下降而营收增长，反映轻资产运营倾向，制约半导体制造重资产投入。
  \item 巴拿马半导体贸易逆差持续扩大，显性比较优势指数始终低于1；STEM毕业生比例偏低，水电基础设施可靠性有待提升，本地供应商和客户生态系统处于早期阶段。
\end{itemize}
\begin{figure}[htbp]
\centering
\includegraphics[width=0.88\linewidth]{figures/fig_150.jpg}
\caption{Figure 2 - Figure 2.2. Geographical distribution of high-tech firms in Panama, 2022 In terms of sectors of economic activity, the majority of high-tech firms were engaged in computer programming, consultancy and related activities (56\% of high-tech firms), followed by th}
\end{figure}
\begin{figure}[htbp]
\centering
\includegraphics[width=0.88\linewidth]{figures/fig_152.jpg}
\caption{Figure 2 - The high-tech ecosystem of Panama is characterised by a large share of micro firms, with around 74\% of the firms having fewer than 10 employees and just 2.0\% of the firms more than 100 employees. This reflects the early-stage nature of this ecosystem and the s}
\end{figure}
\subsection{宏观政策与金融稳定}
\textbf{用于校准利率、通胀、财政约束和金融稳定叙事。}
\begin{itemize}
  \item 2025年11月，IMF发布技术援助报告，总结了为加纳统计服务局（GSS）开发月度宏观环境增长指标（MIEG）的成果。这项由Data for Decisions产品资助的能力建设项目，核心目标是建立一套与《IMF季度国民账户手册（2017）》对齐的月度宏观环境监测框架。报告显示，...
\end{itemize}
\section{结语}
后续需验证：猪周期拐点时间（MS vs 市场预期）、HBM定价弹性（GS vs 彭博一致预期）、美国对中国逆变器限制的实质影响、中国智慧城市增长持续性、以及AI货币化在SaaS领域的拐点信号。
\newpage
\section{覆盖概览}
投行/券商 63 篇；战略咨询 6 篇；智库/国际机构 7 篇。\par
投行覆盖：GS 20篇 / Bernstein 13篇 / MS 11篇 / NOM 8篇 / JPM 4篇 / BARC 2篇 / HSBC 2篇 / UBS 2篇 / Citi 1篇\par
\section{Disclaimer}
{\scriptsize\color{refgray}
This community is a paid learning and information-sharing community created for general educational purposes only. The membership fee is charged solely for access to community discussions, educational content, organization, and operational costs. It is not an advisory fee, management fee, performance fee, brokerage fee, or compensation for personalized financial, investment, legal, tax, or professional advice.\par
I participate and share information solely as an individual and for educational discussion among community members. I am not acting as a financial adviser, investment adviser, broker, dealer, portfolio manager, legal adviser, tax adviser, accountant, fiduciary, or any other licensed professional adviser. Nothing shared in this community constitutes, and should not be interpreted as, financial advice, investment advice, legal advice, tax advice, a recommendation, solicitation, offer, endorsement, instruction, or guarantee to buy, sell, hold, short, trade, or invest in any security, cryptocurrency, fund, derivative, strategy, or other financial product.\par
All content shared in this community, including messages, discussions, links, files, charts, examples, market commentary, watchlists, AI-generated or AI-polished summaries, and third-party materials, is general, impersonal, and educational in nature. It does not take into account any individual's financial situation, investment objectives, risk tolerance, investment horizon, liquidity needs, tax status, legal restrictions, or suitability requirements. No content should be relied upon as suitable for any specific person, account, portfolio, or financial decision.\par
Information shared in this community may be incomplete, inaccurate, outdated, speculative, AI-generated, AI-edited, or based on third-party internet sources that have not been independently verified. I make no representation or warranty as to the accuracy, completeness, timeliness, reliability, or suitability of any information shared.\par
Financial markets involve substantial risk, including the possible loss of principal. Past performance, historical data, hypothetical examples, simulated results, backtests, personal opinions, or educational examples do not guarantee future results. Each member is solely responsible for conducting their own research, exercising independent judgment, and making their own decisions. Before making any financial, investment, legal, or tax decision, members should consult qualified and licensed professionals.\par
Participation in this community, payment of any membership fee, or communication with me or other members does not create any adviser-client, fiduciary, professional, contractual, agency, partnership, employment, or other advisory relationship. By joining or participating in this community, you acknowledge and agree that you are solely responsible for your own decisions, actions, transactions, profits, losses, risks, and consequences, and that I shall not be liable for any direct, indirect, incidental, consequential, financial, legal, tax, or other loss or damage arising from or related to any content, discussion, information, or material shared in this community.\par
}
\newpage
\begin{center}
{\Large 更多详情报告kcdesk.com}\par\vspace{0.8cm}
\includegraphics[width=0.58\linewidth]{\detokenize{../../prompts/zsxq_img.jpg}}
\end{center}
\end{document}
